\documentclass[11pt,twoside]{article}

% ---------------------------------------------------------------------------
% Packages
% ---------------------------------------------------------------------------
\usepackage{amsmath,amssymb}
\usepackage{booktabs}
\usepackage[ruled,vlined,linesnumbered]{algorithm2e}
\usepackage{hyperref}
\usepackage{graphicx}
\usepackage{color}
\usepackage{cite}
\usepackage[margin=1in]{geometry}
\usepackage{enumitem}

% ---------------------------------------------------------------------------
% Theorem environments
% ---------------------------------------------------------------------------
\usepackage{amsthm}
\newtheorem{theorem}{Theorem}[section]
\newtheorem{lemma}[theorem]{Lemma}
\newtheorem{proposition}[theorem]{Proposition}
\newtheorem{corollary}[theorem]{Corollary}
\newtheorem{definition}{Definition}[section]
\newtheorem{assumption}{Assumption}[section]
\newtheorem{remark}{Remark}[section]

% ---------------------------------------------------------------------------
% Math notation
% ---------------------------------------------------------------------------
\newcommand{\R}{\mathbb{R}}
\newcommand{\bA}{\mathbf{A}}
\newcommand{\bb}{\mathbf{b}}
\newcommand{\bc}{\mathbf{c}}
\newcommand{\bx}{\mathbf{x}}
\newcommand{\by}{\mathbf{y}}
\newcommand{\bs}{\mathbf{s}}
\newcommand{\bD}{\mathbf{D}}
\newcommand{\bM}{\mathbf{M}}
\newcommand{\bP}{\mathbf{P}}
\newcommand{\bDelta}{\mathbf{\Delta}}
\newcommand{\bmu}{\boldsymbol{\mu}}
\newcommand{\cO}{\mathcal{O}}
\newcommand{\norm}[1]{\left\lVert#1\right\rVert}
\newcommand{\inner}[1]{\left\langle#1\right\rangle}
\newcommand{\diag}{\operatorname{diag}}
\newcommand{\cond}{\kappa}
\newcommand{\blkdiag}{\operatorname{blkdiag}}

% ---------------------------------------------------------------------------
% Title
% ---------------------------------------------------------------------------
\title{
  Multi-GPU Mixed-Precision Inexact Interior Point Methods \\
  for Large-Scale Dense Linear Programming
}

\author{
  Author One$^{1}$ \and Author Two$^{2}$ \\
  \small $^{1}$Affiliation One \qquad $^{2}$Affiliation Two
}

\date{\today}

\begin{document}
\maketitle

\begin{abstract}
  We propose a multi-GPU mixed-precision inexact interior point method for
  large-scale dense linear programming.  The normal-equations matrix is
  partitioned into diagonal blocks, each factorized on a separate GPU in
  low precision (FP32 or TF32) and used as a block-diagonal preconditioner
  for conjugate gradient or iterative refinement.  The linear-solve
  tolerance is adaptively scheduled as a function of the IPM duality
  measure~$\mu_k$, from $10^{-4}$ in early iterations to $10^{-10}$ near
  optimality, with a guaranteed FP64 fallback.  We prove polynomial
  convergence of the predictor-corrector IPM when the relative residual
  satisfies $\eta_k \le \bar{c}\mu_k^\alpha$, and analyze how inexactness
  propagates through the primal and dual feasibility conditions.
  Experiments on dense instances up to $m=200$, $n=600$ show that
  fixed-accuracy low-precision and block-Jacobi solvers fail to converge,
  while the adaptive solver matches the exact baseline in both iteration
  count and solution quality.  The framework draws on the error-management
  philosophy of quantum interior point methods, recasting it as a practical
  strategy for heterogeneous GPU computing.
\end{abstract}

% ===================================================================
\section{Introduction}
\label{sec:intro}
% ===================================================================

Linear programming (LP) is a fundamental technology in operations research,
machine learning, quantitative finance, and engineering design.  Modern
applications routinely produce LP instances with tens of thousands of
variables and constraints: portfolio optimization over large asset
universes, sparse recovery and compressed sensing in signal processing,
optimal transport in machine learning, and network flow problems in
logistics and supply chain management, to name a few.  When the constraint
matrix is dense---as it is in many of these applications after
preprocessing, basis transformations, or when the problem is naturally
fully coupled---the computational bottleneck shifts decisively to the
linear algebra at each iteration of the optimization algorithm.

Interior point methods (IPMs)~\cite{Wright1997,Nesterov1994} are among the
most reliable and theoretically well-understood algorithms for solving LPs.
A primal-dual IPM follows the central path toward optimality by solving, at
each iteration, a Newton system that, after elimination of the primal and
dual slack variables, reduces to the $m \times m$ symmetric positive
definite \emph{normal equations}
\begin{equation}
  \bM_k \bDelta\by_k = \br_k, \qquad
  \bM_k := \bA \bD_k^2 \bA^T,
  \label{eq:normal}
\end{equation}
where $\bA \in \R^{m \times n}$ is the constraint matrix,
$\bD_k = \diag(\sqrt{x_1/s_1},\dots,\sqrt{x_n/s_n})$ encodes the
primal-dual scaling at iteration $k$, and $\br_k$ is the assembled
right-hand side of the Mehrotra predictor-corrector
scheme~\cite{Mehrotra1992}.  The matrix $\bM_k$ must be ``solved''---in
the sense of computing $\bDelta\by_k$ that satisfies
\eqref{eq:normal}---at every IPM iteration, and this is by far the most
expensive operation in the algorithm.  In the codebase, $\bM_k$ is formed
by the function \texttt{compute\_normal\_equation} in
\texttt{lp\_problem.py} as \texttt{A\_scaled = A * D\_diag[None, :]}
followed by \texttt{A\_scaled @ A\_scaled.T}, exactly implementing
$\bM_k = (\bA\bD_k)(\bA\bD_k)^T$.

For large-scale \emph{dense} LPs, where $m \approx 10^3$--$10^5$ and
$\bA$ has no exploitable sparsity pattern, the standard approach is to
factor $\bM_k$ via dense Cholesky decomposition $\bM_k = \bL\bL^T$,
which costs $\cO(m^3)$ floating-point operations and must be repeated at
every single IPM iteration (typically 15--40 iterations for a well-tuned
implementation).  The total cost of the Cholesky factorizations alone can
easily reach $10^{12}$--$10^{15}$ operations, dominating the solution
time.  For instance, at $m = 10^4$, each Cholesky factorization involves
approximately $1.7 \times 10^{11}$ operations; across 20 IPM iterations,
this amounts to $3.4 \times 10^{12}$ operations, requiring minutes to
hours on a single CPU core.

Two recent developments, one in hardware and one in theory, motivate a
fundamentally different approach to this bottleneck.

\paragraph{Hardware trend: the GPU precision gap.}
Modern GPU architectures deliver extraordinarily high throughput for
dense linear algebra operations (GEMM, SYRK, TRSM) through tensor
cores~\cite{NVIDIA2023}.  However, the throughput is strongly
precision-dependent: on the NVIDIA H100 GPU, FP64 (double precision)
delivers approximately 34 TFLOPS, while FP32 delivers 67 TFLOPS and TF32
(tensor-float-32) reaches 990 TFLOPS---a nearly $30\times$ gap between
the highest-throughput low-precision mode and the FP64 mode required for
classical IPM linear algebra.  This precision gap means that a Cholesky
factorization performed entirely in FP64 on a GPU leaves most of the
hardware's computational capacity unused.  If we could perform the bulk
of the factorization in FP32 or TF32 and recover accuracy through a small
number of correction steps in FP64, the effective throughput could
increase by an order of magnitude or more.

\paragraph{Algorithmic trend: inexact IPM theory.}
Independently, the theory of \emph{inexact} interior point
methods~\cite{Bellavia1998,Monteiro2003} has established that the Newton
direction need not be computed to machine precision at every IPM
iteration.  It suffices that the relative residual of the linear solve,
\begin{equation}
  \eta_k := \frac{\norm{\bM_k\widehat{\bDelta\by}_k - \br_k}_2}
                   {\norm{\br_k}_2},
  \label{eq:eta-def-intro}
\end{equation}
be bounded by a quantity that shrinks with the IPM's progress toward
optimality: $\eta_k \le \bar{c}\,\mu_k^\alpha$, where $\mu_k =
\bx_k^T\bs_k / n$ is the current duality measure.  In plain terms, the
linear solve can be quite inaccurate when the IPM is far from optimality
($\mu_k$ large), but must become progressively more accurate as the
solver converges ($\mu_k \to 0$).  This theoretical insight transforms
the linear-solve accuracy from a \emph{fixed} requirement (machine
precision at every iteration) into a \emph{dynamic} resource that can be
allocated proportionally to the IPM's needs.  The residual $\eta_k$ is
computed in the code by \texttt{\_compute\_eta} in
\texttt{linear\_solver.py} as $\texttt{np.linalg.norm(M @ delta\_y -
r) / np.linalg.norm(r)}$, always evaluated in FP64 regardless of the
internal precision used by the solver.

Our central observation is that these two developments are
\emph{complementary}.  The GPU precision gap makes low-precision
factorization fast; inexact IPM theory makes low-accuracy solves
acceptable.  What is needed is a bridge between them: a
precision-and-accuracy scheduling mechanism that decides, at each IPM
iteration, \emph{how much} accuracy is required of the linear solve and
\emph{in what precision} to compute the factorization, given the current
duality measure $\mu_k$.

We propose such a bridge in this paper: the \emph{multi-GPU mixed-precision
inexact interior point method} for large-scale dense linear programming.
The algorithm combines three ideas.  First, the normal-equations matrix
$\bM_k$ is partitioned into $p$ diagonal blocks, each assigned to a
distinct GPU and factorized independently via dense Cholesky in low
precision (FP32 or TF32), yielding a block-diagonal preconditioner
$\bP_k = \blkdiag(\bM_{11},\dots,\bM_{pp})$.  This block-partition strategy
is implemented in \texttt{\_build\_block\_diagonal\_preconditioner} in
\texttt{linear\_solver.py}, which slices $\bM_k$ into \texttt{num\_blocks}
contiguous diagonal blocks and factors each block via Cholesky.  Second, the
Newton step is computed by preconditioned conjugate gradient or mixed-precision
iterative refinement, reusing the low-precision block factors to correct the
solution---a \emph{factorize-once, solve-many} pattern that amortizes the
$\cO(m^3)$ factorization cost across the predictor step, the corrector step,
and all refinement corrections.  Third, and most critically, the target
accuracy for the linear solve is adaptively scheduled as a piecewise function
of $\mu_k$, implemented by \texttt{\_mu\_to\_eta\_target} in
\texttt{linear\_solver.py}:
$\eta_{\text{target}}(\mu) = 10^{-4}$ when $\mu > 10^{-2}$ (early IPM, loose
accuracy), $\eta_{\text{target}}(\mu) = 10^{-7}$ when $10^{-6} < \mu \le
10^{-2}$ (mid IPM), and $\eta_{\text{target}}(\mu) = 10^{-10}$ when $\mu \le
10^{-6}$ (late IPM, near-exact solves).  A fallback to FP64 Cholesky
guarantees convergence when the low-precision factorization cannot meet the
target.

The approach is inspired, in part, by recent work on quantum interior point
methods (QIPM)~\cite{Kerenidis2020}.  In QIPM, a quantum linear system
algorithm (QLSA) provides an approximate Newton direction whose error is
bounded in $\ell_2$-norm, and the QIPM literature has developed a detailed
theory of how to schedule QLSA accuracy across IPM iterations.  Our setting
is strictly classical---the error source is floating-point round-off and
block-diagonal approximation rather than quantum sampling noise---but the
mathematical structure is isomorphic: both frameworks must cope with an
inexact linear solve at each iteration, and both rely on the outer IPM's
ability to absorb bounded linear-system error.  The adaptive tolerance
schedule developed here can be viewed as a classical analogue of the QLSA
precision schedule in QIPM, and our analysis of feasibility drift
(Section~\ref{sec:error-prop}) parallels the feasibility-preserving
Newton reformulations studied in the QIPM literature.

The paper makes the following contributions.  It presents a convergence
analysis of the Mehrotra predictor-corrector IPM with inexact Newton
directions governed by the condition $\eta_k \le \bar{c}\mu_k^\alpha$,
showing that polynomial complexity $\cO(\sqrt{n}\log(1/\varepsilon))$ is
retained (Section~\ref{sec:theory}, Theorem~\ref{th:conv}).  It develops
a multi-GPU block-diagonal Cholesky preconditioner that trades off-diagonal
coupling for embarrassingly parallel factorization across GPUs, analyzes its
approximation quality through the generalized eigenvalue
bound~\eqref{eq:cond-preconditioned}, and presents the
\texttt{LowPrecisionCholeskySolver} and \texttt{BlockJacobiPCGSolver}
algorithms (Section~\ref{sec:precond}, Algorithms~\ref{alg:lp-cholesky}
and~\ref{alg:bj-pcg}).  It designs and implements a $\mu$-driven adaptive
refinement solver that dynamically adjusts the linear-solve accuracy target
across three phases, mathematically encoding this as the piecewise
function~\eqref{eq:mu-schedule} and embedding it in an iterative refinement
loop with a guaranteed FP64 fallback (Section~\ref{sec:adapt},
Algorithm~\ref{alg:adapt}).  And it provides numerical experiments on dense
LP instances up to $m=200$, $n=600$ that reveal a sharp empirical finding:
fixed-accuracy low-precision and block-Jacobi solvers \emph{fail to converge
on every problem size tested}, while the adaptive solver tracks the exact
baseline in both iteration count and solution quality
(Section~\ref{sec:experiments}, Table~\ref{tab:full-results}).

The remainder of the paper is organized as follows.  Section~\ref{sec:prelim}
reviews the primal-dual IPM framework, the normal-equations reduction, the
Mehrotra predictor-corrector algorithm, and the fraction-to-boundary step
rule, establishing precise notation matched to the reference implementation.
Section~\ref{sec:theory} develops the inexact IPM theory, defining the
relative residual $\eta_k$, analyzing its propagation through the IPM step
(primal feasibility, dual feasibility, and duality measure), stating the
polynomial convergence theorem, and discussing the feasible vs.\ infeasible
distinction.  Section~\ref{sec:precond} presents the block-partition
structure of $\bM_k$, the $2\times 2$ block Cholesky with Schur complement,
the block-diagonal preconditioner, the three-level precision strategy, and
the \texttt{LowPrecisionCholeskySolver} and
\texttt{BlockJacobiPCGSolver} algorithms.  Section~\ref{sec:adapt}
introduces the $\mu$-driven tolerance schedule, the complete
\texttt{AdaptiveRefinementSolver} algorithm, a detailed design analysis,
and the connection to QIPM.  Section~\ref{sec:experiments} reports the
numerical results.  Section~\ref{sec:conclusion} concludes with a summary
and directions for future work.

% ===================================================================
\section{Preliminaries}
\label{sec:prelim}
% ===================================================================

\subsection{Standard-form linear programming}

We consider the primal linear program in standard form and its Lagrangian dual:
\begin{align}
  \min_{\bx \in \R^n}\;& \bc^T \bx
    \quad \text{subject to}\quad
    \bA\bx = \bb,\;\; \bx \ge 0,
    \label{eq:primal} \\[4pt]
  \max_{\by \in \R^m,\; \bs \in \R^n}\;& \bb^T \by
    \quad \text{subject to}\quad
    \bA^T\by + \bs = \bc,\;\; \bs \ge 0.
    \label{eq:dual}
\end{align}
Here $\bA \in \R^{m \times n}$ is the constraint matrix (typically
$m \le n$), $\bb \in \R^m$ is the right-hand side vector, and
$\bc \in \R^n$ is the cost vector.  The component-wise inequalities
$\bx \ge 0$ and $\bs \ge 0$ define the non-negative orthant.
Problem~\eqref{eq:primal} is called the \emph{primal} LP and
Problem~\eqref{eq:dual} is the \emph{dual} LP; the variable $\bx$ is the
primal solution, $\by$ is the dual solution (or Lagrange multiplier for
the equality constraint $\bA\bx = \bb$), and $\bs$ is the dual slack
variable.

The first-order necessary and sufficient optimality conditions for the pair
\eqref{eq:primal}--\eqref{eq:dual} are the Karush--Kuhn--Tucker (KKT)
system
\begin{equation}
  \begin{aligned}
    \bA\bx &= \bb, \\
    \bA^T\by + \bs &= \bc, \\
    \bx \circ \bs &= \mathbf{0}, \\
    \bx &\ge \mathbf{0},\quad \bs \ge \mathbf{0},
  \end{aligned}
  \label{eq:kkt}
\end{equation}
where $\bx \circ \bs \in \R^n$ denotes the Hadamard (entrywise) product,
i.e., $(\bx \circ \bs)_i = x_i s_i$ for $i = 1,\dots,n$, and
$\mathbf{0}$ is the vector of all zeros.  The first equation of~\eqref{eq:kkt}
is primal feasibility; the second is dual feasibility; and the third is the
complementarity condition, which requires $x_i s_i = 0$ for every $i$.  Because $\bx \ge 0$ and $\bs \ge 0$, the
complementarity condition is equivalent to the scalar equation
$\bx^T\bs = 0$, which states that the duality gap is zero at optimality.

\subsection{The central path and the duality measure}

Primal-dual interior point methods replace the non-smooth
complementarity condition of~\eqref{eq:kkt} with the \emph{perturbed}
condition
\begin{equation}
  \bx \circ \bs = \mu \mathbf{1}_n,
  \label{eq:perturbed-compl}
\end{equation}
where $\mu > 0$ is a centering parameter and $\mathbf{1}_n$ is the vector
of all ones in $\R^n$.  For each fixed $\mu > 0$, the perturbed KKT system
\begin{equation}
  \bA\bx = \bb,\qquad
  \bA^T\by + \bs = \bc,\qquad
  \bx \circ \bs = \mu \mathbf{1}_n,\qquad
  \bx > \mathbf{0},\; \bs > \mathbf{0},
  \label{eq:perturbed-kkt}
\end{equation}
has a unique solution $(\bx(\mu), \by(\mu), \bs(\mu))$ under mild
assumptions (e.g., $\bA$ has full row rank and a strictly feasible
interior point exists).  The set
\begin{equation}
  \mathcal{C} := \bigl\{ (\bx(\mu), \by(\mu), \bs(\mu)) \mid \mu > 0 \bigr\}
\end{equation}
is called the \emph{central path}.  As $\mu \to 0$, the central path
converges to a solution of the original KKT system~\eqref{eq:kkt}.

Taking the inner product of both sides of~\eqref{eq:perturbed-compl} with
$\mathbf{1}_n$ yields
\begin{equation}
  \bx^T\bs = \sum_{i=1}^n x_i s_i = n\mu,
  \label{eq:duality-gap}
\end{equation}
so that the quantity
\begin{equation}
  \mu := \frac{\bx^T\bs}{n}
  \label{eq:mu-def}
\end{equation}
is the average complementarity product at the current iterate.  We refer
to $\mu$ as the \emph{duality measure}; it equals the normalized duality
gap and is the natural measure of progress toward optimality in a
primal-dual IPM.  In the code, $\mu$ is computed exactly as in
\eqref{eq:mu-def}: $\texttt{mu} = \texttt{float}(\texttt{np.dot}(x, s) /
n)$.

\subsection{Newton system and the normal equations}

At each iteration of a primal-dual IPM, we linearize the perturbed KKT
system~\eqref{eq:perturbed-kkt} around the current iterate
$(\bx_k, \by_k, \bs_k)$ with $\bx_k > 0$, $\bs_k > 0$, and compute a
search direction $(\bDelta\bx_k, \bDelta\by_k, \bDelta\bs_k)$.
Linearizing~\eqref{eq:perturbed-compl} yields the third equation
\begin{equation}
  \bS_k \bDelta\bx_k + \bX_k \bDelta\bs_k = \sigma_k \mu_k \mathbf{1}_n
  - \bX_k \bs_k,
  \label{eq:linearized-compl}
\end{equation}
where $\bX_k := \diag(\bx_k)$, $\bS_k := \diag(\bs_k)$, $\mu_k :=
\bx_k^T\bs_k / n$, and $\sigma_k \in [0,1]$ is the \emph{centering
parameter}.  Setting $\sigma_k = 1$ gives a pure centering direction
(toward the point $(\bx(\mu_k), \by(\mu_k), \bs(\mu_k))$ on the central
path), while $\sigma_k = 0$ gives the pure affine-scaling direction
(toward optimality, ignoring centrality).  The full linearized system is
the $(2n+m) \times (2n+m)$ \emph{augmented system}
\begin{equation}
  \begin{bmatrix}
    \mathbf{0} & \bA^T & \mathbf{I} \\
    \bA        & \mathbf{0} & \mathbf{0} \\
    \mathbf{0} & \mathbf{0} & \bS_k
  \end{bmatrix}
  \begin{bmatrix}
    \bDelta\bx_k \\ \bDelta\by_k \\ \bDelta\bs_k
  \end{bmatrix}
  =
  \begin{bmatrix}
    \bc - \bA^T\by_k - \bs_k \\
    \bb - \bA\bx_k \\
    \sigma_k \mu_k \mathbf{1}_n - \bX_k\bS_k \mathbf{1}_n
  \end{bmatrix}.
  \label{eq:augmented-system}
\end{equation}

Eliminating $\bDelta\bs_k$ via the third block row gives
$\bDelta\bs_k = \sigma_k \mu_k \bX_k^{-1}\mathbf{1}_n - \bs_k -
\bX_k^{-1}\bS_k\bDelta\bx_k$.  Substituting into the first two block rows
and eliminating $\bDelta\bx_k$ reduces the system to the $m \times m$
\emph{normal equations}
\begin{equation}
  \boxed{\bM_k \bDelta\by_k = \br_k}, \qquad
  \bM_k := \bA \bD_k^2 \bA^T,
  \label{eq:normal-detailed}
\end{equation}
where the \emph{scaling matrix} is
\begin{equation}
  \bD_k := \diag\bigl(d_1^{(k)},\dots,d_n^{(k)}\bigr),\qquad
  d_i^{(k)} := \sqrt{\frac{x_i^{(k)}}{s_i^{(k)}}},\quad i=1,\dots,n,
  \label{eq:scaling}
\end{equation}
and the right-hand side $\br_k \in \R^m$ is assembled from the residuals
and the centering term, as detailed in Section~\ref{sec:mehrotra} below.
In the codebase, $\bM_k$ is formed by the function
$\texttt{compute\_normal\_equation}$, which computes
\texttt{A\_scaled = A * D\_diag[None, :]} followed by
\texttt{A\_scaled @ A\_scaled.T}.

The normal-equations matrix $\bM_k$ is symmetric positive definite (SPD)
whenever $\bx_k > 0$, $\bs_k > 0$, and $\bA$ has full row rank.  This
SPD property is the key advantage of the normal-equations formulation: it
allows the use of Cholesky factorization, conjugate gradient, and other
methods specialized for SPD systems.

Once $\bDelta\by_k$ has been computed (exactly or approximately, as
discussed in Section~\ref{sec:theory}), the primal and dual slack
directions are recovered via
\begin{align}
  \bDelta\bx_k &=
    -\bx_k
    + \frac{\sigma_k\mu_k}{\bs_k}
    - \frac{\bx_k}{\bs_k} \circ \bigl(\bc - \bA^T\by_k - \bs_k\bigr)
    + \frac{\bx_k}{\bs_k} \circ (\bA^T\bDelta\by_k)
    - \frac{\bDelta\bx_k^{\text{aff}} \circ \bDelta\bs_k^{\text{aff}}}{\bs_k},
    \label{eq:recover-dx} \\[4pt]
  \bDelta\bs_k &=
    \bigl(\bc - \bA^T\by_k - \bs_k\bigr) - \bA^T\bDelta\by_k,
    \label{eq:recover-ds}
\end{align}
where all divisions and multiplications involving vectors are
component-wise, and the final term in~\eqref{eq:recover-dx} is the
second-order Mehrotra correction (see Section~\ref{sec:mehrotra}).

\subsection{The Mehrotra predictor-corrector algorithm}
\label{sec:mehrotra}

We adopt the widely used Mehrotra~\cite{Mehrotra1992} predictor-corrector
scheme, which is the workhorse of virtually all practical IPM
implementations.  The algorithm computes, at each iteration, two
directions using the \emph{same} coefficient matrix $\bM_k$ but different
right-hand sides.

\subsubsection{Step 1: Predictor (affine-scaling) direction}

Set $\sigma = 0$ (no centering) and omit the second-order correction.
Define the dual residual at the current iterate
\begin{equation}
  \br_d^{(k)} := \bc - \bA^T\by_k - \bs_k.
  \label{eq:rd}
\end{equation}
The predictor right-hand side (in $\R^m$) is
\begin{equation}
  \br_k^{\text{aff}} := \bb + \bA\Bigl(\frac{\bx_k}{\bs_k}
                         \circ \br_d^{(k)}\Bigr).
  \label{eq:r-aff}
\end{equation}
Solve the normal equations (exactly or inexactly) to obtain
$\bDelta\by_k^{\text{aff}}$, and recover the affine primal and dual slack
directions $\bDelta\bx_k^{\text{aff}}$ and $\bDelta\bs_k^{\text{aff}}$
via~\eqref{eq:recover-dx}--\eqref{eq:recover-ds} with $\sigma=0$ and no
correction term.

\subsubsection{Step 2: Fraction-to-boundary step lengths}

The step length to the boundary of the non-negative orthant is determined
by the \emph{fraction-to-boundary} rule.  For a current iterate $\bv > 0$
and a direction $\bDelta\bv$, the maximum step that preserves
non-negativity is
\begin{equation}
  \alpha_{\max}(\bv, \bDelta\bv) :=
    \min\left\{
      -\frac{v_i}{\Delta v_i}
      \;\middle|\;
      i \in \{1,\dots,n\},\; \Delta v_i < 0
    \right\},
  \label{eq:alpha-max}
\end{equation}
with $\alpha_{\max} := \tau$ if no component of $\bDelta\bv$ is negative.
To keep the iterates strictly positive and bounded away from the boundary
(which would cause the scaling matrix $\bD_k$ to blow up), the actual step
length is taken as
\begin{equation}
  \boxed{\alpha(\bv, \bDelta\bv) :=
    \min\bigl\{\tau \cdot \alpha_{\max}(\bv, \bDelta\bv),\; \tau\bigr\}},
  \label{eq:step-length}
\end{equation}
where $\tau \in (0,1)$ is the fraction-to-boundary parameter (typically
$\tau = 0.99$ or $0.9995$).  The step length for the dual variable $\by$
is taken equal to that of the dual slack $\bs$, since $\by$ is
unconstrained in sign.  These step lengths are computed for the affine
direction to obtain $\alpha_p^{\text{aff}} :=
\alpha(\bx_k, \bDelta\bx_k^{\text{aff}})$ and $\alpha_d^{\text{aff}} :=
\alpha(\bs_k, \bDelta\bs_k^{\text{aff}})$.

\subsubsection{Step 3: Centering parameter}

Let
\begin{equation}
  \mu_k^{\text{aff}} :=
    \frac{1}{n}
      \bigl(\bx_k + \alpha_p^{\text{aff}}\bDelta\bx_k^{\text{aff}}\bigr)^T
      \bigl(\bs_k + \alpha_d^{\text{aff}}\bDelta\bs_k^{\text{aff}}\bigr)
  \label{eq:mu-aff}
\end{equation}
be the duality measure that would result from a full step along the affine
direction using the computed step lengths.  The ratio $\mu_k^{\text{aff}}
/ \mu_k$ measures how much progress the pure affine direction makes.
Mehrotra's heuristic sets the centering parameter adaptively as
\begin{equation}
  \boxed{\sigma_k := \left(\frac{\mu_k^{\text{aff}}}{\mu_k}\right)^3},
  \label{eq:sigma}
\end{equation}
with the safeguard $\sigma_k := 0$ when $\mu_k \le 10^{-16}$.  The
exponent $3$ is standard practice: when the affine direction makes
substantial progress ($\mu_k^{\text{aff}} \ll \mu_k$), $\sigma_k$ becomes
very small, biasing the corrector step toward optimality; conversely,
when the affine direction makes little progress, $\sigma_k \approx 1$ and
the corrector step emphasizes centrality.

\subsubsection{Step 4: Corrector (combined) direction}

The corrector right-hand side assembles three components:
\begin{align}
  \br_k^{\text{corr},1} &:= \bb
    + \bA\Bigl(\frac{\bx_k}{\bs_k} \circ \br_d^{(k)}\Bigr),
    \label{eq:corr1} \\[4pt]
  \br_k^{\text{corr},2} &:= \sigma_k\mu_k\,
    \bA\Bigl(\frac{\mathbf{1}_n}{\bs_k}\Bigr),
    \label{eq:corr2} \\[4pt]
  \br_k^{\text{corr},3} &:=
    \bA\Bigl(\frac{\bDelta\bx_k^{\text{aff}}
                 \circ \bDelta\bs_k^{\text{aff}}}{\bs_k}\Bigr),
    \label{eq:corr3}
\end{align}
where \eqref{eq:corr1} is the linearized residual term (identical in form
to the predictor RHS~\eqref{eq:r-aff}), \eqref{eq:corr2} is the centering
term (which vanishes when $\sigma_k \to 0$), and \eqref{eq:corr3} is
Mehrotra's second-order nonlinear correction that uses the
Hadamard product of the affine directions to approximate the quadratic
term in the Taylor expansion of the complementarity condition.  The
full corrector right-hand side is
\begin{equation}
  \br_k^{\text{corr}} :=
    \br_k^{\text{corr},1}
    - \br_k^{\text{corr},2}
    + \br_k^{\text{corr},3}.
  \label{eq:corr-rhs}
\end{equation}
The normal equations $\bM_k\bDelta\by_k = \br_k^{\text{corr}}$ are solved
with the \emph{same} coefficient matrix $\bM_k$ as the predictor step
(which is critical for efficiency: one factorization serves two solves).
The primal and dual directions are then recovered via
\eqref{eq:recover-dx}--\eqref{eq:recover-ds} with $\sigma = \sigma_k$ and
the correction term included.

\subsubsection{Step 5: Step length and update}

The step lengths for the combined (corrector) direction are computed via
\eqref{eq:step-length}:
\begin{equation}
  \alpha_p^{(k)} := \alpha(\bx_k, \bDelta\bx_k),\qquad
  \alpha_d^{(k)} := \alpha(\bs_k, \bDelta\bs_k).
  \label{eq:step-lengths}
\end{equation}
The iterate is updated as
\begin{equation}
  \begin{aligned}
    \bx_{k+1} &:= \bx_k + \alpha_p^{(k)} \bDelta\bx_k, \\
    \bs_{k+1} &:= \bs_k + \alpha_d^{(k)} \bDelta\bs_k, \\
    \by_{k+1} &:= \by_k + \alpha_d^{(k)} \bDelta\by_k.
  \end{aligned}
  \label{eq:update}
\end{equation}
Notice that the dual variable $\by$ steps with $\alpha_d^{(k)}$ (the same
step length as the dual slack), which is standard practice since both
live in the dual space and using the smaller of $\alpha_p$ and $\alpha_d$
for $\by$ would be overly conservative.

\subsubsection{Convergence monitoring}

At each iteration, the relative primal and dual residuals are computed as
\begin{align}
  \rho_p^{(k)} &:= \frac{\norm{\bb - \bA\bx_k}_2}
                        {1 + \norm{\bb}_2},
    \label{eq:primal-res} \\[4pt]
  \rho_d^{(k)} &:= \frac{\norm{\bc - \bA^T\by_k - \bs_k}_2}
                        {1 + \norm{\bc}_2},
    \label{eq:dual-res}
\end{align}
where the $1 + \norm{\cdot}_2$ denominators follow standard practice for
scale-invariant stopping criteria~\cite{Wright1997}.  The IPM is
considered converged when
\begin{equation}
  \mu_k < \varepsilon,\quad
  \rho_p^{(k)} < \varepsilon,\quad
  \rho_d^{(k)} < \varepsilon,
  \label{eq:ipm-conv-test}
\end{equation}
with a typical tolerance of $\varepsilon = 10^{-8}$.

\subsection{Conditioning of the normal equations}

A fundamental numerical difficulty with the normal-equations formulation
is that the spectral condition number of $\bM_k$ can be dramatically worse
than that of the original constraint matrix $\bA$.  From the definition
$\bM_k = \bA\bD_k^2\bA^T$ and the submultiplicativity of the spectral
norm,
\begin{equation}
  \cond(\bM_k) = \cond(\bA\bD_k^2\bA^T)
    \le \cond(\bA)^2 \cdot
    \Bigl(\frac{\max_i d_i^{(k)}}{\min_i d_i^{(k)}}\Bigr)^2,
  \label{eq:cond-bound}
\end{equation}
where $\cond(\cdot) := \sigma_{\max}(\cdot) / \sigma_{\min}(\cdot)$ is the
ratio of largest to smallest singular value.  The factor $(\max d_i /
\min d_i)^2$ reflects the spread of the diagonal scaling entries.  In the
early iterations of an IPM, the entries of $\bx_k$ and $\bs_k$ are
typically of comparable magnitude, so $d_i^{(k)} = \sqrt{x_i/s_i}$ is
moderate and $\bM_k$ is well-conditioned.  However, as the IPM converges
and $\mu_k \to 0$, the primal and dual variables separate: some $x_i$
approach zero (active constraints) while the corresponding $s_i$ remain
bounded away from zero (or vice versa).  Consequently, the ratio
$\max_i d_i / \min_i d_i$ can span many orders of magnitude—often $10^8$
or larger in the final iterations—and $\cond(\bM_k)$ can exceed
$10^{16}$, the reciprocal of machine epsilon in FP64.

This progressive ill-conditioning is the central numerical challenge that
any low-precision or inexact linear solver must confront inside an IPM.
It motivates the adaptive precision schedule developed in
Section~\ref{sec:adapt} and the fallback mechanisms embedded in the
solvers (see Section~\ref{sec:precond}).

% ===================================================================
\section{Inexact IPM Theory}
\label{sec:theory}
% ===================================================================

In an exact IPM, the normal equations~\eqref{eq:normal-detailed} are solved
to machine precision at each iteration.  In an \emph{inexact} IPM, we permit
an approximate solution, trading linear-solver accuracy for computational
speed.  This section formalizes the notion of inexactness, analyzes how the
linear-solve error propagates through the IPM step, and states the
convergence guarantees that underpin the adaptive scheduling strategy of
Section~\ref{sec:adapt}.

\subsection{Inexact Newton directions}

Let $\widehat{\bDelta\by}_k$ denote the \emph{computed} (approximate)
solution of the normal equations at iteration $k$.  We write
\begin{equation}
  \bM_k \widehat{\bDelta\by}_k = \br_k + \be_k^{\text{lin}},
  \label{eq:inexact-system}
\end{equation}
where
\begin{equation}
  \be_k^{\text{lin}} := \bM_k \widehat{\bDelta\by}_k - \br_k
  \label{eq:e-lin}
\end{equation}
is the \emph{linear-system residual}.  The exact Newton direction
$\bDelta\by_k^*$ satisfies $\bM_k\bDelta\by_k^* = \br_k$, so
$\be_k^{\text{lin}} = \bM_k(\widehat{\bDelta\by}_k - \bDelta\by_k^*)$.

\begin{definition}[Relative residual]
  \label{def:eta}
  The \emph{relative residual} of the linear solve is
  \begin{equation}
    \boxed{\eta_k :=
      \frac{\norm{\be_k^{\text{lin}}}_2}{\norm{\br_k}_2}
      \equiv
      \frac{\norm{\bM_k\widehat{\bDelta\by}_k - \br_k}_2}
           {\norm{\br_k}_2}}.
    \label{eq:eta-def}
  \end{equation}
  When $\norm{\br_k}_2 < 10^{-15}$ we set $\eta_k := 0$.
\end{definition}

This is precisely the quantity tracked by the $\texttt{\_compute\_eta}$
method in the codebase: the residual is evaluated in FP64 regardless of the
internal precision used to compute $\widehat{\bDelta\by}_k$, so $\eta_k$
gives a faithful, precision-independent measure of backward error.

The central question of inexact IPM theory is: \emph{how large can
$\eta_k$ be at iteration $k$ while still guaranteeing that the outer IPM
converges, and at what asymptotic rate?}

\subsection{Error propagation through the IPM step}
\label{sec:error-prop}

\subsubsection{Error in the primal-dual directions}

Given the approximate dual direction $\widehat{\bDelta\by}_k$, the primal
and dual slack directions are recovered via~\eqref{eq:recover-dx}--
\eqref{eq:recover-ds} (omitting the corrector terms for clarity; the
analysis generalizes straightforwardly):
\begin{align}
  \widehat{\bDelta\bx}_k &=
    -\bx_k
    + \frac{\sigma_k\mu_k}{\bs_k}
    - \frac{\bx_k}{\bs_k} \circ \br_d^{(k)}
    + \frac{\bx_k}{\bs_k} \circ (\bA^T\widehat{\bDelta\by}_k),
    \label{eq:inexact-recover-dx} \\[4pt]
  \widehat{\bDelta\bs}_k &=
    \br_d^{(k)} - \bA^T\widehat{\bDelta\by}_k,
    \label{eq:inexact-recover-ds}
\end{align}
where $\br_d^{(k)} = \bc - \bA^T\by_k - \bs_k$ is the dual residual at
the current iterate.  Let $\bDelta\bx_k^*$ and $\bDelta\bs_k^*$ be the
exact directions that would be obtained by using $\bDelta\by_k^*$ in the
same formulas.  Subtracting the exact from the inexact recovery yields the
\emph{direction errors}:
\begin{align}
  \widehat{\bDelta\bx}_k - \bDelta\bx_k^* &=
    \bX_k^{1/2}\bS_k^{-1/2} \circ
    \bA^T(\widehat{\bDelta\by}_k - \bDelta\by_k^*)
    = \bD_k \circ \bA^T(\widehat{\bDelta\by}_k - \bDelta\by_k^*),
    \label{eq:dx-error} \\[4pt]
  \widehat{\bDelta\bs}_k - \bDelta\bs_k^* &=
    -\bA^T(\widehat{\bDelta\by}_k - \bDelta\by_k^*).
    \label{eq:ds-error}
\end{align}
The error in the dual direction propagates to the primal direction
\emph{scaled by the factor $\bD_k$}, which amplifies the error precisely
when the scaling matrix is ill-conditioned (i.e., in late IPM iterations).

\subsubsection{Propagation to primal feasibility}

The primal residual after taking a step of length $\alpha$ in the inexact
direction is
\begin{align}
  \br_p^{(k+1)}
    &= \bb - \bA(\bx_k + \alpha_k \widehat{\bDelta\bx}_k) \nonumber \\
    &= \bb - \bA\bx_k
       - \alpha_k \bA\bDelta\bx_k^*
       - \alpha_k \bA(\widehat{\bDelta\bx}_k - \bDelta\bx_k^*) \nonumber \\
    &= (1 - \alpha_k)(\bb - \bA\bx_k)
       + \alpha_k \bA\bD_k^2\bA^T
         (\widehat{\bDelta\by}_k - \bDelta\by_k^*),
    \label{eq:feas-prop}
\end{align}
where we used $\bA\bDelta\bx_k^* = \bb - \bA\bx_k$ (the exact direction
restores primal feasibility exactly in one full step) and
\eqref{eq:dx-error} to express the direction error in terms of
$\widehat{\bDelta\by}_k - \bDelta\by_k^*$.  Recalling that
$\bM_k = \bA\bD_k^2\bA^T$ and $\be_k^{\text{lin}} =
\bM_k(\widehat{\bDelta\by}_k - \bDelta\by_k^*)$, we obtain the compact
relation
\begin{equation}
  \boxed{\br_p^{(k+1)} =
    (1 - \alpha_k)\br_p^{(k)} + \alpha_k \be_k^{\text{lin}}}.
  \label{eq:feas-prop-final}
\end{equation}

Equation~\eqref{eq:feas-prop-final} reveals a fundamental tension: the
step length $\alpha_k$ simultaneously reduces the current primal residual
(the first term) and injects the linear-solve error (the second term).
If $\norm{\be_k^{\text{lin}}}$ is large relative to $\norm{\br_p^{(k)}}$,
the primal residual may \emph{grow} rather than shrink.  This phenomenon is
observed in practice: our implementation detects the condition
\begin{equation}
  \mu_k < 10^{-6} \;\text{and}\; \rho_p^{(k)} > 10^{-3}
  \label{eq:feas-loss}
\end{equation}
as an indicator that the linear-solve error has irreparably damaged primal
feasibility, and the IPM is terminated.

\subsubsection{Propagation to dual feasibility}

A similar analysis yields the dual residual after the step:
\begin{equation}
  \br_d^{(k+1)} =
    (1 - \alpha_k)\br_d^{(k)}
    + \alpha_k \bA^T(\widehat{\bDelta\by}_k - \bDelta\by_k^*).
  \label{eq:dual-prop}
\end{equation}
Comparing with~\eqref{eq:feas-prop-final}, the dual residual receives the
\emph{unscaled} direction error, while the primal residual receives it
premultiplied by $\bA\bD_k^2\bA^T = \bM_k$.  Since $\bM_k$ can be highly
ill-conditioned (see Section~\ref{sec:prelim}), the primal residual is
typically more sensitive to inexactness than the dual residual.

\subsubsection{Propagation to the duality measure}

The duality measure after the step is
\begin{equation}
  \mu_{k+1} = \frac{1}{n}
    (\bx_k + \alpha_k\widehat{\bDelta\bx}_k)^T
    (\bs_k + \alpha_k\widehat{\bDelta\bs}_k).
  \label{eq:mu-prop}
\end{equation}
Expanding and using $\bx_k^T\bs_k = n\mu_k$, one can show that
\begin{equation}
  \mu_{k+1} = (1 - \alpha_k(1 - \sigma_k))\mu_k
    + \cO(\alpha_k^2\norm{\cdot}^2)
    + \cO(\alpha_k \eta_k \norm{\br_k}),
  \label{eq:mu-prop-expanded}
\end{equation}
where the final $\cO(\alpha_k \eta_k \norm{\br_k})$ term captures the
effect of the inexactness.  This term must be dominated by the reduction
in $\mu$ for convergence to be maintained.

\subsection{Convergence theorem}
\label{sec:conv-theorem}

The convergence of inexact IPMs is governed by how the allowable residual
$\eta_k$ scales with the duality measure $\mu_k$.  We state the central
result (without full proof; see~\cite{Bellavia1998,Monteiro2003} for
detailed analyses).

\begin{assumption}[Interior starting point]
  \label{as:feas}
  The initial iterate $(\bx_0, \by_0, \bs_0)$ is strictly feasible
  ($\bx_0 > 0$, $\bs_0 > 0$) with $\bA\bx_0 = \bb$ and
  $\bA^T\by_0 + \bs_0 = \bc$, and lies in a suitable neighborhood
  $\cN_2(\beta)$ of the central path:
  \begin{equation}
    \norm{\bx_0 \circ \bs_0 - \mu_0\mathbf{1}_n}_2 \le \beta\mu_0,
    \qquad \beta \in (0,1).
    \label{eq:nbhd}
  \end{equation}
\end{assumption}

\begin{assumption}[Full rank]
  \label{as:rank}
  The constraint matrix $\bA \in \R^{m \times n}$ has full row rank, so
  that $\bM_k = \bA\bD_k^2\bA^T$ is strictly positive definite for all
  $\bx_k,\bs_k > 0$.
\end{assumption}

\begin{theorem}[Polynomial convergence of inexact IPM]
  \label{th:conv}
  Let Assumptions~\ref{as:feas} and~\ref{as:rank} hold.  Suppose that at
  each iteration $k$ of the Mehrotra predictor-corrector IPM, the linear
  solve satisfies
  \begin{equation}
    \boxed{\eta_k \le \bar{c}\,\mu_k^\alpha},
    \label{eq:eta-condition}
  \end{equation}
  for constants $\bar{c} > 0$ and $\alpha \in (0,1]$ independent of $k$.
  Then there exists a constant $\delta > 0$ such that
  \begin{equation}
    \mu_{k+1} \le \Bigl(1 - \frac{\delta}{\sqrt{n}}\Bigr)\mu_k,
    \label{eq:mu-linear}
  \end{equation}
  and the method converges to a point with $\mu \le \varepsilon$ in at most
  \begin{equation}
    K = \cO\bigl(\sqrt{n}\,\log(1/\varepsilon)\bigr)
    \label{eq:iter-bound}
  \end{equation}
  iterations.
\end{theorem}

\begin{proof}[Proof sketch]
  The proof adapts the standard analysis of polynomial IPM convergence
  (see, e.g.,~\cite{Wright1997, Nesterov1994}) by tracking the effect of
  the inexactness residual through the predictor and corrector steps
  separately.

  \textbf{Predictor step.}  The affine direction with $\sigma=0$ reduces
  $\mu$ by a factor that, in the exact case, is $(1 - \delta/\sqrt{n})$.
  With inexactness, the reduction is attenuated by an $\cO(\eta_k)$ term.
  Condition~\eqref{eq:eta-condition} ensures that $\eta_k \le \bar{c}\mu_k$
  (when $\alpha \ge 1$) or $\eta_k \le \bar{c}\mu_k^\alpha$ (when
  $\alpha < 1$), so the inexactness budget is always small relative to the
  current duality gap, preserving the polynomial reduction rate.

  \textbf{Corrector step.}  The centering component (with $\sigma_k$ given
  by~\eqref{eq:sigma}) restores proximity to the central path.  The
  inexactness adds a perturbation that is $\cO(\eta_k^2)$ in the centrality
  measure~\eqref{eq:nbhd}, which is dominated by the linear term when
  $\eta_k$ is sufficiently small.

  \textbf{Asymptotic behavior.}  When $\mu_k \ll 1$ and $\alpha$ is chosen
  appropriately, $\eta_k$ shrinks rapidly with $\mu_k$.  If
  $\alpha \ge 1$, then $\eta_k$ shrinks at least linearly with $\mu_k$,
  and the quadratic convergence of Mehrotra's method (in the final
  iterations) is preserved.  If $\alpha < 1$, the inexactness decays more
  slowly than $\mu_k$, and the asymptotic rate degrades to linear, though
  the polynomial complexity bound~\eqref{eq:iter-bound} still holds.

  \textbf{Feasibility drift.}  Equation~\eqref{eq:feas-prop-final} shows
  that the primal residual evolves as $\br_p^{(k+1)} = (1-\alpha_k)\br_p^{(k)}
  + \alpha_k \be_k^{\text{lin}}$.  Under~\eqref{eq:eta-condition}, the
  injected error $\be_k^{\text{lin}}$ is bounded by
  $\bar{c}\mu_k^\alpha\norm{\br_k}$, which is $\cO(\mu_k^\alpha)$ when the
  RHS $\br_k$ is of moderate size.  Since $\mu_k$ itself decreases
  geometrically by~\eqref{eq:mu-linear}, the cumulative feasibility drift
  over $K$ iterations is bounded, and the final primal infeasibility is
  $\cO(\varepsilon^\alpha)$.
\end{proof}

\begin{remark}
  The condition $\eta_k \le \bar{c}\mu_k^\alpha$ subsumes several special
  cases from the literature.  Bellavia~\cite{Bellavia1998} uses
  $\eta_k \le \gamma\mu_k$ (i.e., $\alpha=1$, $\bar{c}=\gamma$), while
  Monteiro and O'Neal~\cite{Monteiro2003} consider $\eta_k \le
  \gamma\min\{\mu_k,1\}$, which is equivalent to $\alpha=1$ with a cap.
  Our adaptive schedule (Section~\ref{sec:adapt}) uses a piecewise-constant
  approximation to a $\mu^\alpha$ envelope, which is easier to implement
  and tune in practice.
\end{remark}

\subsection{Feasible vs.\ infeasible inexact IPM}
\label{sec:feas-vs-infeas}

Equation~\eqref{eq:feas-prop-final} highlights a critical distinction
between two classes of inexact IPM formulations:

\begin{definition}[Inexact feasible IPM]
  An IPM is called \emph{inexact feasible} if the Newton direction is
  constructed so that primal and dual feasibility are preserved exactly,
  even when the linear solve is approximate.  This requires that the
  inexactness residual $\be_k^{\text{lin}}$ lie in the null space of
  $\bA$ (for primal feasibility) or that the direction recovery
  formulas be modified to project out the error~\cite{Kerenidis2020}.
\end{definition}

\begin{definition}[Inexact infeasible IPM]
  An IPM is called \emph{inexact infeasible} if no special mechanism is
  used to preserve feasibility, and the residuals evolve according to
  \eqref{eq:feas-prop-final}--\eqref{eq:dual-prop}.  This is the natural
  setting for the normal-equations formulation with an approximate
  $\widehat{\bDelta\by}_k$.
\end{definition}

The prototype implemented in this paper is an \emph{inexact infeasible}
IPM.  The practical implications depend on the phase of the solve.  In the
early iterations, $\mu_k$ is large, the normal-equations matrix is
well-conditioned, and the linear-solve error $\eta_k$ is small (either
because the low-precision factorization is accurate, or because iterative
refinement converges quickly); the feasibility
drift~\eqref{eq:feas-prop-final} is negligible.  In the mid iterations, as
$\mu_k$ decreases, $\kappa(\bM_k)$ grows and the low-precision
factorization loses accuracy; the adaptive schedule tightens
$\eta_{\text{target}}$, triggering more refinement iterations to keep
$\eta_k$ within bounds; primal feasibility may degrade moderately, but the
IPM corrector step (with $\sigma_k > 0$) restores centrality and helps
contain the drift.  In the late iterations, when $\mu_k < 10^{-6}$, even
small linear-solve errors can translate into large primal residuals because
$\bD_k$ amplifies them via~\eqref{eq:dx-error}; at this stage, the
adaptive schedule demands near-exact solves ($\eta_{\text{target}} =
10^{-10}$), and the fallback to FP64 Cholesky ensures that the final
solution accuracy is not compromised.  Finally, the code monitors the
condition~\eqref{eq:feas-loss} and terminates early if the inexactness has
pushed the IPM into an irrecoverable infeasible state: this is a pragmatic
safeguard that avoids wasting computation on iterates that can no longer
converge and provides diagnostic feedback for tuning the tolerance
schedule.

The inexact infeasible formulation is both simpler to implement and more
natural for the normal-equations framework than a feasibility-preserving
alternative (e.g., via null-space projection or augmented KKT systems).
The price is the need for the heuristic safeguards described above; the
numerical results in Section~\ref{sec:experiments} confirm that these
safeguards are effective in practice.

\begin{remark}[Connection to QIPM]
  Quantum interior point methods~\cite{Kerenidis2020} face an analogous
  trade-off: the quantum linear system algorithm (QLSA) outputs an
  approximate solution whose error must be absorbed by the outer IPM.
  The QIPM literature has therefore developed a detailed theory of
  inexact \emph{feasible} IPMs, where the Newton reformulation (often via
  null-space methods) preserves feasibility regardless of linear-solve
  accuracy.  Our work can be viewed as a classical analogue: we use
  low-precision GPU arithmetic in place of QLSA noise, and iterative
  refinement in place of quantum state tomography, but the mathematical
  structure of inexactness is similar.  This connection suggests that
  feasibility-preserving reformulations developed for QIPM could also
  benefit classical GPU-accelerated IPMs, and vice versa.
\end{remark}

% ===================================================================
\section{Multi-GPU Block Cholesky Preconditioner}
\label{sec:precond}
% ===================================================================

The normal-equations matrix $\bM_k = \bA\bD_k^2\bA^T$ is an $m \times m$
symmetric positive definite matrix.  For large $m$ (e.g., $m \ge 10^4$),
a full dense Cholesky factorization $\bM_k = \bL\bL^T$ costs $\cO(m^3)$
operations and requires $\cO(m^2)$ storage, which may exceed the memory of
a single GPU.  This section develops a \emph{block-diagonal Cholesky
preconditioner} that distributes the factorization across multiple GPUs
and serves as the computational core of all the inexact solvers described
in this paper.

\subsection{Block partition of the normal-equations matrix}

We partition the index set $\{1,\dots,m\}$ into $p$ contiguous blocks of
sizes $m_1,\dots,m_p$ (with $\sum_{i=1}^p m_i = m$).  The normal-equations
matrix inherits the corresponding block structure:
\begin{equation}
  \bM_k =
  \begin{bmatrix}
    \bM_{11} & \bM_{12} & \cdots & \bM_{1p} \\[4pt]
    \bM_{21} & \bM_{22} & \cdots & \bM_{2p} \\[4pt]
    \vdots   & \vdots   & \ddots & \vdots   \\[4pt]
    \bM_{p1} & \bM_{p2} & \cdots & \bM_{pp}
  \end{bmatrix},
  \qquad
  \bM_{ij} \in \R^{m_i \times m_j},
  \label{eq:block-partition}
\end{equation}
where $\bM_{ij} = \bA_i\bD_k^2\bA_j^T$ and $\bA_i \in \R^{m_i \times n}$ is
the $i$-th row-block of the constraint matrix.  Because $\bM_k$ is symmetric,
$\bM_{ji} = \bM_{ij}^T$.

In our multi-GPU setting, the $i$-th diagonal block $\bM_{ii}$ (and the
corresponding row-block $\bA_i$) is assigned to GPU $i$.  The block sizes
$m_i$ are chosen to balance GPU memory utilization; a simple equal-split
strategy sets $m_i \approx m/p$ for all $i$, as implemented in
\texttt{\_build\_block\_diagonal\_preconditioner} by the slicing
\texttt{start = i * m // p; end = (i+1) * m // p}.

\subsection{Block Cholesky factorization with Schur complement}

To understand the relationship between a full Cholesky factorization and
the block-diagonal approximation, it is instructive to examine the $p=2$
case in detail.  Write
\begin{equation}
  \bM_k =
  \begin{bmatrix}
    \bM_{11} & \bM_{12} \\
    \bM_{21} & \bM_{22}
  \end{bmatrix},
  \qquad
  \bM_{11} \in \R^{m_1 \times m_1},\;
  \bM_{22} \in \R^{m_2 \times m_2},
  \label{eq:2x2-block}
\end{equation}
with $\bM_{12}^T = \bM_{21}$.  The \emph{block Cholesky factorization}
(generalizing the standard block LDL$^T$ decomposition to the SPD case) is
\begin{equation}
  \bM_k =
  \begin{bmatrix}
    \bL_{11} & \mathbf{0} \\
    \bM_{21}\bL_{11}^{-T} & \bL_{22}
  \end{bmatrix}
  \begin{bmatrix}
    \bL_{11}^T & \bL_{11}^{-1}\bM_{12} \\
    \mathbf{0} & \bL_{22}^T
  \end{bmatrix},
  \label{eq:block-cholesky}
\end{equation}
where
\begin{align}
  \bM_{11} &= \bL_{11}\bL_{11}^T,
    \label{eq:chol-M11} \\[4pt]
  \bS_{22} &:= \bM_{22} - \bM_{21}\bM_{11}^{-1}\bM_{12}
    = \bL_{22}\bL_{22}^T.
    \label{eq:schur}
\end{align}
The matrix $\bS_{22}$ is the \emph{Schur complement} of $\bM_{11}$ in
$\bM_k$.  It is SPD whenever $\bM_k$ is SPD.  Equation~\eqref{eq:schur}
reveals the essential difficulty of distributed block factorization: even
if diagonal blocks $\bM_{11}$ and $\bM_{22}$ are factorized independently on
separate GPUs, the second block's factor $\bL_{22}$ requires the Schur
complement update $\bS_{22} = \bM_{22} - \bM_{21}\bM_{11}^{-1}\bM_{12}$,
which couples the two GPUs through the off-diagonal block $\bM_{21}$.
Computing $\bM_{21}\bM_{11}^{-1}\bM_{12}$ involves solving triangular
systems with $\bL_{11}$ and then a matrix-matrix multiplication, both of
which require communication when $\bM_{11}$ and $\bM_{22}$ reside on
different devices.

For $p > 2$, the block Cholesky proceeds recursively: after factorizing
$\bM_{11} = \bL_{11}\bL_{11}^T$, the trailing $(p-1) \times (p-1)$
submatrix is updated by the Schur complement of $\bM_{11}$, and the
process repeats.  The communication pattern is all-to-all in the worst
case, and the total communication volume grows with $p$.

\subsection{Block-diagonal preconditioner: trading exactness for parallelism}

Rather than attempting a full distributed block Cholesky (which would
require substantial inter-GPU communication for the Schur complement
updates), we construct a \emph{block-diagonal preconditioner}:
\begin{equation}
  \boxed{\bP_k := \blkdiag(\bM_{11}, \bM_{22}, \dots, \bM_{pp})},
  \label{eq:precond-def}
\end{equation}
where each diagonal block is factorized \emph{independently}:
\begin{equation}
  \bM_{ii} = \widetilde{\bL}_i \widetilde{\bL}_i^T,
  \qquad i = 1,\dots,p.
  \label{eq:local-cholesky}
\end{equation}
The tildes on $\widetilde{\bL}_i$ indicate that the factorization may be
performed in reduced precision (see Section~\ref{sec:precision-strategy}).

The preconditioner $\bP_k$ is an \emph{approximation} to $\bM_k$: it
coincides with $\bM_k$ on the diagonal blocks but ignores the off-diagonal
coupling.  The quality of the approximation is characterized by the
generalized eigenvalue problem: if $\lambda$ is an eigenvalue of
$\bP_k^{-1}\bM_k$, then the condition number of the preconditioned system
governs the convergence rate of preconditioned Krylov methods.  A standard
bound (via block Gershgorin) is
\begin{equation}
  \cond\bigl(\bP_k^{-1/2}\bM_k\bP_k^{-1/2}\bigr)
    \le 1 + \max_{1 \le i \le p}\;
    \sum_{\substack{j=1 \\ j \ne i}}^{p}
    \bigl\|\bM_{ii}^{-1/2}\bM_{ij}\bM_{jj}^{-1/2}\bigr\|_2,
  \label{eq:cond-preconditioned}
\end{equation}
where the sum over $j \ne i$ quantifies the \emph{off-diagonal coupling
strength}.  When the off-diagonal blocks are small in norm relative to the
diagonal blocks, the preconditioned operator is close to the identity and
Krylov methods converge in few iterations.  Conversely, when the couplings
are strong (e.g., when $\bA$ is dense with no natural block structure),
the preconditioner is weaker and more inner iterations are required.

This design reflects a deliberate trade-off: we sacrifice the exactness of
a full Cholesky solve for \emph{embarrassingly parallel} factorization
across $p$ GPUs with zero inter-GPU communication during the factorization
phase.  The resulting error is absorbed by the inexact IPM framework of
Section~\ref{sec:theory} and, when necessary, corrected by iterative
refinement or Krylov acceleration.

\subsection{Precision strategy}
\label{sec:precision-strategy}

Modern GPUs exhibit dramatically different throughput across floating-point
formats.  For instance, the NVIDIA H100 GPU delivers approximately 34 TFLOPS
in FP64, 67 TFLOPS in FP32, and up to 990 TFLOPS in TF32 (tensor-core)
mode~\cite{NVIDIA2023}.  To exploit this hierarchy while protecting the
convergence guarantees of the outer IPM, we adopt a \emph{three-level}
precision strategy.

At the first level, \textbf{factorization precision (FP32 / TF32)}, each
diagonal block $\bM_{ii}$ is downcast to FP32 (or TF32, which shares the
FP32 exponent range with reduced mantissa precision) and factorized via
dense Cholesky:
\begin{equation}
  \bM_{ii}^{\text{lo}} = \texttt{float32}(\bM_{ii}),\qquad
  \bM_{ii}^{\text{lo}} \approx \widetilde{\bL}_i \widetilde{\bL}_i^T.
  \label{eq:fp32-factor}
\end{equation}
The downcast uses the function \texttt{\_safe\_cast\_f32} which clips
values exceeding the FP32 range ($\approx 3.4 \times 10^{38}$) to prevent
overflow.  A small diagonal regularization $\rho \approx 10^{-10}$ is added
before factorization to guard against singularity due to round-off.  If
the Cholesky factorization fails (e.g., due to an indefinite matrix in
FP32), the code falls back to an eigendecomposition
$\bM_{ii}^{\text{lo}} = \bV \bLambda \bV^T$ with eigenvalues clamped to a
minimum of $10^{-12}$.

At the second level, \textbf{residual computation (FP64)}, regardless of
the precision used for factorization and solve, the linear-system residual
is always evaluated in full FP64:
\begin{equation}
  \be_k^{\text{lin}} = \br_k - \bM_k\widehat{\bDelta\by}_k
  \quad\text{(computed in float64)}.
  \label{eq:fp64-residual}
\end{equation}
This is critical: it provides a precision-independent backward error
measure $\eta_k$ that faithfully reflects the true linear-solve accuracy,
preventing a situation where the solver believes it has converged but the
direction is actually corrupted by low-precision round-off.

At the third level, \textbf{iterate update (FP64)}, the Newton direction
$\widehat{\bDelta\by}_k$ and the IPM iterates $(\bx_{k+1}, \by_{k+1},
\bs_{k+1})$ are accumulated in FP64.  This ensures that the outer
optimization is not limited by the precision of the linear algebra; only
the per-iteration factorization benefits from hardware-accelerated
low-precision arithmetic.

This three-level hierarchy is a direct adaptation of the mixed-precision
iterative refinement framework of Carson and Higham~\cite{Carson2018} to
the structured linear systems arising inside an IPM.  The key difference is
that our precision schedule is not static but adapts to the IPM's progress
through the $\mu$-driven tolerance schedule of Section~\ref{sec:adapt}.

\subsection{Solver 1: LowPrecisionCholeskySolver}
\label{sec:lp-cholesky}

The simplest inexact solver in our framework is the
\texttt{LowPrecisionCholeskySolver}, which applies the three-level
precision strategy to the \emph{full} normal-equations matrix without
block partitioning (i.e., $p=1$).  It simulates a single-GPU tensor-core
Cholesky solve and serves as a minimal baseline for evaluating the effect
of reduced precision alone.

\begin{algorithm}[t]
  \caption{\texttt{LowPrecisionCholeskySolver.solve(M, r)}}
  \label{alg:lp-cholesky}
  \KwIn{SPD matrix $\bM \in \R^{m \times m}$, right-hand side $\br \in \R^m$}
  \KwOut{Solution $\widehat{\bDelta\by}$, relative residual $\eta$}

  \tcp{Step 1: Downcast to FP32 with overflow guard}
  $\widetilde{\bM} \gets \texttt{\_safe\_cast\_f32}(\bM)$ \tcp*{clip to $[-10^{37}, 10^{37}]$, cast to float32}
  $\widetilde{\br} \gets \texttt{\_safe\_cast\_f32}(\br)$\;

  \tcp{Step 2: Diagonal regularization}
  $\widetilde{\bM}.\text{diag} \,\texttt{+=}\;\texttt{float32}(10^{-10})$\;

  \tcp{Step 3: Factorization and solve in FP32}
  \Try{}{
    $(\widetilde{\bL}, \texttt{lower}) \gets \texttt{cho\_factor}(\widetilde{\bM})$\;
    $\widehat{\bDelta\by}^{\text{lo}} \gets \texttt{cho\_solve}(\widetilde{\bL}, \widetilde{\br})$\;
  }\Catch{LinAlgError}{
    $(\bw, \bV) \gets \texttt{eigh}(\widetilde{\bM})$ \tcp*{fallback: eigendecomposition}
    $\bw \gets \max(\bw, 10^{-12})$\;
    $\widehat{\bDelta\by}^{\text{lo}} \gets \bV\bigl((\bV^T\widetilde{\br}) \oslash \bw\bigr)$\;
  }

  \tcp{Step 4: Upcast to FP64}
  $\widehat{\bDelta\by} \gets \texttt{float64}(\widehat{\bDelta\by}^{\text{lo}})$\;

  \tcp{Step 5: Fallback if FP32 produced NaN}
  \If{$\text{any}(\texttt{isnan}(\widehat{\bDelta\by}))$}{
    $\widehat{\bDelta\by} \gets \texttt{solve}(\bM, \br)$ \tcp*{FP64 direct solve}
  }

  \tcp{Step 6: Compute relative residual in FP64}
  $\eta \gets \norm{\bM\widehat{\bDelta\by} - \br}_2 \;/\; \norm{\br}_2$\;
  \Return $(\widehat{\bDelta\by}, \eta)$\;
\end{algorithm}

Algorithm~\ref{alg:lp-cholesky} reproduces the exact logic of the
implementation.  Several design choices merit comment.  The overflow guard
(\texttt{\_safe\_cast\_f32}) clips entries to $\pm 10^{37}$ to prevent Inf
values in FP32, which has a maximum representable value of approximately
$3.4 \times 10^{38}$; for well-scaled IPM iterates this guard is rarely
triggered, but it is essential for robustness with random problem data.
The eigendecomposition fallback handles the rare case where round-off in
FP32 causes $\widetilde{\bM}$ to lose positive definiteness: rather than
aborting, we compute $\widetilde{\bM} = \bV\diag(\bw)\bV^T$ and solve via
$\widehat{\bDelta\by}^{\text{lo}} = \bV\diag(1/\bw)\bV^T \widetilde{\br}$,
which is always well-defined with the eigenvalue clamping.  The NaN
fallback (Step~5 of the algorithm) catches catastrophic precision loss: if
the FP32 computation produced non-finite values, we solve the original FP64
system directly; this is rarely executed in practice but ensures the solver
never silently returns garbage.  Finally, the residual is always evaluated
in FP64 (Step~6, matching \texttt{\_compute\_eta}) so that $\eta$
accurately reflects the backward error, not the forward error of the
low-precision solve.

\subsection{Solver 2: BlockJacobiPCGSolver}
\label{sec:bj-pcg}

The \texttt{BlockJacobiPCGSolver} extends the block-diagonal
preconditioner~\eqref{eq:precond-def} to a full Krylov solver.  It is the
primary solver designed for the multi-GPU setting envisioned in this
paper.  The code is in \texttt{BlockJacobiPCGSolver.solve} and
\texttt{\_build\_block\_diagonal\_preconditioner}.

\subsubsection{Preconditioner construction}

Given the SPD matrix $\bM_k$ and the number of blocks $p$, the
preconditioner is built as follows.  The index range $[0, m)$ is partitioned
into $p$ contiguous slices $[\texttt{start}_i, \texttt{end}_i)$ of
approximately equal size.  For each slice $i$, the diagonal block
$\bM_{ii} = \bM_k[\texttt{start}_i{:}\texttt{end}_i,\,
\texttt{start}_i{:}\texttt{end}_i]$ is extracted, and a small diagonal shift
$\bM_{ii} \gets \bM_{ii} + 10^{-12}\, \mathbf{I}$ is applied for numerical
stability.  The Cholesky factorization $\bM_{ii} =
\widetilde{\bL}_i\widetilde{\bL}_i^T$ is computed via
\texttt{scipy.linalg.cho\_factor}, and the factor $\widetilde{\bL}_i$ and
the slice boundaries are stored for subsequent solves.

The preconditioner operator $P = \texttt{LinearOperator}(m, m,
\text{matvec}=\dots)$ applies $\bP_k^{-1}\bv$ to any vector
$\bv \in \R^m$ by solving, for each block $i$ independently,
\begin{equation}
  \widetilde{\bL}_i \widetilde{\bL}_i^T \bz_i = \bv_i,
  \label{eq:block-solve}
\end{equation}
where $\bv_i = \bv[\texttt{start}_i{:}\texttt{end}_i]$.  The triangular
solves are executed sequentially on the CPU in the reference
implementation; on a multi-GPU system, each block solve runs on its
assigned GPU concurrently.

\subsubsection{Conjugate gradient solve}

\begin{algorithm}[t]
  \caption{\texttt{BlockJacobiPCGSolver.solve(M, r)}}
  \label{alg:bj-pcg}
  \KwIn{SPD matrix $\bM \in \R^{m \times m}$, right-hand side $\br \in \R^m$}
  \KwOut{Solution $\widehat{\bDelta\by}$, relative residual $\eta$}

  \tcp{Step 1: Build block-diagonal preconditioner}
  $p \gets \min(\texttt{num\_blocks}, m)$\;
  \For{$i \in \{0,\dots,p-1\}$}{
    $\texttt{start} \gets i \cdot m \;/\; p$,
    $\texttt{end} \gets (i+1) \cdot m \;/\; p$\;
    $\bB \gets \bM[\texttt{start}{:}\texttt{end},\; \texttt{start}{:}\texttt{end}].
      \texttt{copy}()$\;
    $\bB.\text{diag} \,\texttt{+=}\; 10^{-12}$\;
    $\widetilde{\bL}_i \gets \texttt{cho\_factor}(\bB)$\;
  }
  $\bP^{-1} \gets \texttt{LinearOperator}(x \mapsto \text{apply } \widetilde{\bL}_i
    \text{ to each slice of } x)$\;

  \tcp{Step 2: Preconditioned conjugate gradient}
  $\widehat{\bDelta\by} \gets \texttt{cg}(
    \bM,\; \br,\;
    M = \bP^{-1},\;
    \texttt{rtol} = \texttt{cg\_tol},\;
    \texttt{maxiter} = \texttt{max\_iter}
  )$\;

  \tcp{Step 3: Residual in FP64}
  $\eta \gets \norm{\bM\widehat{\bDelta\by} - \br}_2 \;/\; \norm{\br}_2$\;
  \Return $(\widehat{\bDelta\by}, \eta)$\;
\end{algorithm}

Several properties of this solver are worth noting.  First, the parameter
$p$ provides a continuum of inexactness that controls the
accuracy--parallelism trade-off.  When $p = 1$, $\bP_k = \bM_k$ and PCG
converges in one iteration (the preconditioner is the exact solve),
degenerating to a direct method with no parallelism.  When $p = m$,
$\bP_k = \diag(\bM_k)$ is the Jacobi (diagonal) preconditioner; each CG
iteration is extremely cheap ($\cO(m)$ with no triangular solves), but
convergence may require many iterations on ill-conditioned problems.  For
intermediate $p$ (e.g., $p = 4$ or $8$), each GPU holds a block of size
$m/p$, and the preconditioner quality depends on the off-diagonal coupling
as bounded by~\eqref{eq:cond-preconditioned}.  Second, the CG solver
terminates when the preconditioned residual norm falls below
\texttt{cg\_tol} (default $10^{-6}$); because the true residual $\eta$ is
computed separately in FP64, the CG tolerance can be set looser than the
IPM's accuracy requirement, and the adaptive refinement step
(Section~\ref{sec:adapt}) can compensate for remaining error.  Third, each
diagonal block is factored independently, so the factorization phase is
trivially parallel across GPUs with zero communication; communication
occurs only during the CG iteration in the matrix-vector product $\bM_k
\bv$, which is the standard communication pattern of distributed CG and is
well supported by MPI and NCCL collectives.  Fourth, since the predictor
and corrector steps use the same coefficient matrix $\bM_k$ but different
right-hand sides, the block-diagonal factorization $\{\widetilde{\bL}_i\}$
computed once can be reused for both solves; furthermore, the iterative
refinement corrections (Algorithm~\ref{alg:adapt}) also reuse the same
factors, making this a factorize-once, solve-many scheme that is
particularly efficient on GPU hardware.

% ===================================================================
\section{Adaptive Refinement Strategy}
\label{sec:adapt}
% ===================================================================

The solvers described in Section~\ref{sec:precond} all operate with a fixed
accuracy target: the \texttt{ExactCholeskySolver} solves to machine
precision, the \texttt{LowPrecisionCholeskySolver} accepts whatever
accuracy FP32 provides, and the \texttt{BlockJacobiPCGSolver} terminates
when the preconditioned CG residual meets a user-specified tolerance.  This
section introduces the \texttt{AdaptiveRefinementSolver}, which
\emph{dynamically} adjusts the linear-solve accuracy target as a function
of the IPM's current duality measure $\mu_k$.  The core idea is that the
linear solve need not be equally accurate at every IPM iteration: early
iterations can tolerate large errors, while late iterations demand
near-exact solves.

\subsection{Motivation: why adaptivity matters}

Consider the trajectory of a typical IPM solve.  In the first few
iterations, $\mu_k$ is large (often $\mu_0 \sim 10^0$ to $10^2$ for the
initial point $\bx_0 = \bs_0 = \mathbf{1}$), the scaling matrix
$\bD_k$ is well-conditioned (all $d_i \approx 1$), and the normal-equations
matrix $\bM_k = \bA\bD_k^2\bA^T$ has a condition number close to that of
$\bA\bA^T$.  A low-precision Cholesky factorization in FP32 produces a
direction accurate to $\eta_k \approx 10^{-5}$ to $10^{-4}$, which is more
than sufficient when $\mu_k \sim 10^0$; by Theorem~\ref{th:conv}, the
inexactness budget is $\bar{c}\mu_k^\alpha \approx \bar{c} \cdot 1^\alpha
= \bar{c}$, easily satisfied.  Insisting on $\eta_k \le 10^{-10}$ in these
early iterations would waste computation with no benefit to convergence.

As the solve progresses and $\mu_k$ decreases through $10^{-3}, 10^{-6},
10^{-9}$, two things happen simultaneously: the convergence theory demands a
tighter $\eta_k$ to maintain the polynomial reduction
rate~\eqref{eq:mu-linear}, and the conditioning of $\bM_k$ deteriorates
(recall \eqref{eq:cond-bound}), making it harder to achieve a given
$\eta_k$ with the same low-precision factorization.  The adaptive strategy
reconciles these opposing trends: it demands higher accuracy precisely when
the problem requires it, while relaxing accuracy when the problem is
forgiving.

\subsection{The $\mu$-driven tolerance schedule}
\label{sec:mu-schedule}

The mapping from duality measure to linear-solve tolerance is implemented
by the function \texttt{\_mu\_to\_eta\_target} in
\texttt{linear\_solver.py}.  We formalize it as follows.

\begin{definition}[$\mu$-driven tolerance schedule]
  \label{def:mu-schedule}
  The \emph{target relative residual} for the linear solve at an IPM
  iteration with duality measure $\mu > 0$ is
  \begin{equation}
    \boxed{\eta_{\text{target}}(\mu) :=
    \begin{cases}
      10^{-4},  & \mu > 10^{-2}
        \quad \text{(Phase I: early IPM — loose)}, \\[6pt]
      10^{-7},  & 10^{-6} < \mu \le 10^{-2}
        \quad \text{(Phase II: mid IPM — moderate)}, \\[6pt]
      10^{-10}, & \mu \le 10^{-6}
        \quad \text{(Phase III: late IPM — near-exact)}.
    \end{cases}}
    \label{eq:mu-schedule}
  \end{equation}
  When $\mu$ is \texttt{None} (i.e., the solver is called outside an IPM
  context), $\eta_{\text{target}} := 10^{-4}$.
\end{definition}

The three phases correspond to distinct numerical regimes:

\paragraph{Phase I: $\mu > 10^{-2}$ (early IPM).}
The duality gap is large, the central path is broad, and the IPM is far
from the ill-conditioned region near optimality.  The linear solve can be
quite inexact ($\eta_{\text{target}} = 10^{-4}$) without endangering
convergence.  In practice, a single FP32 Cholesky solve typically achieves
$\eta_k \approx 10^{-5}$ to $10^{-4}$ on well-conditioned $\bM_k$, so
\emph{no refinement iterations are needed} in this phase.  This is where
the adaptive strategy realizes its largest computational savings.

\paragraph{Phase II: $10^{-6} < \mu \le 10^{-2}$ (mid IPM).}
The duality gap has decreased by several orders of magnitude.  The scaling
matrix $\bD_k$ has developed some spread, and $\cond(\bM_k)$ has
increased.  The target tightens to $\eta_{\text{target}} = 10^{-7}$,
which a single FP32 Cholesky solve can no longer reliably achieve.  One
or two refinement iterations are typically needed to close the gap between
the FP32 initial solve and the target.

\paragraph{Phase III: $\mu \le 10^{-6}$ (late IPM).}
The IPM is approaching optimality.  The normal-equations matrix may be
severely ill-conditioned ($\cond(\bM_k) > 10^{12}$).  The target
$\eta_{\text{target}} = 10^{-10}$ demands near-exact solves.  Multiple
refinement iterations are usually required; if the refinement loop exhausts
its budget without meeting the target, the solver falls back to a full FP64
Cholesky solve, which is guaranteed to achieve machine-precision accuracy
(assuming $\bM_k$ is not numerically singular).

\subsection{Relationship to the convergence theory}

The piecewise-constant schedule~\eqref{eq:mu-schedule} can be viewed as a
practical approximation to the theoretical condition $\eta_k \le
\bar{c}\mu_k^\alpha$ from Theorem~\ref{th:conv}.  Choosing
$\bar{c} = 10^{-4}$ and $\alpha = \log_{10}(10^{-4}/10^{-10}) /
\log_{10}(10^2/10^{-6}) \approx 0.75$, the envelope
$10^{-4}\mu_k^{0.75}$ passes through the points
$(\mu=10^{-2}, \eta=10^{-5.5})$, $(\mu=10^{-6}, \eta=10^{-8.5})$,
$(\mu=10^{-10}, \eta=10^{-11.5})$, which the step function
\eqref{eq:mu-schedule} bounds from above at every $\mu$.  Thus the
schedule is a \emph{conservative} approximation to the theoretical
requirement: it demands somewhat higher accuracy than strictly necessary
at each $\mu$, providing a margin of safety against the approximations in
the convergence analysis (e.g., the proof's reliance on worst-case
constants).

\subsection{The AdaptiveRefinementSolver}
\label{sec:adapt-algorithm}

Algorithm~\ref{alg:adapt} gives the complete pseudocode for the
\texttt{AdaptiveRefinementSolver.solve} method, matching the implementation
in \texttt{linear\_solver.py} lines~355--472.  The algorithm combines a
low-precision initial solve (FP32 Cholesky) with up to
\texttt{max\_refinement\_steps} rounds of iterative refinement, all
governed by the $\mu$-driven tolerance schedule.

\begin{algorithm}[t]
  \caption{\texttt{AdaptiveRefinementSolver.solve(M, r, mu)}}
  \label{alg:adapt}
  \KwIn{SPD matrix $\bM \in \R^{m \times m}$, right-hand side $\br \in \R^m$,
        duality measure $\mu \in \R_+ \cup \{\texttt{None}\}$}
  \KwOut{Solution $\widehat{\bDelta\by} \in \R^m$, relative residual $\eta$}
  \BlankLine

  \tcp{Phase 0: Determine target tolerance}
  $\eta_{\text{target}} \gets \textsc{MuToEtaTarget}(\mu)$
    \tcp*{Definition~\ref{def:mu-schedule}}
  $\texttt{num\_ref} \gets 0$,
  $\texttt{used\_fallback} \gets \texttt{False}$\;

  \tcp{Phase 1: Low-precision factorization and initial solve}
  $\widetilde{\bM} \gets \texttt{\_safe\_cast\_dtype}(\bM,\; \texttt{float32})$\;
  $\widetilde{\br} \gets \texttt{\_safe\_cast\_dtype}(\br,\; \texttt{float32})$\;
  $\widetilde{\bM}.\text{diag} \,\texttt{+=}\; 10^{-10}$ \tcp*{regularize}

  \Try{}{
    $(\widetilde{\bL}, \texttt{lower}) \gets \texttt{cho\_factor}(\widetilde{\bM})$\;
    $\widehat{\bDelta\by}^{\text{lo}} \gets \texttt{cho\_solve}(\widetilde{\bL}, \widetilde{\br})$\;
    $(\texttt{factor}, \texttt{use\_eigh}) \gets
      ((\widetilde{\bL}, \texttt{lower}),\; \texttt{False})$\;
  }\Catch{LinAlgError}{
    $(\bw, \bV) \gets \texttt{eigh}(\widetilde{\bM})$\;
    $\bw \gets \max(\bw, 10^{-12})$\;
    $\widehat{\bDelta\by}^{\text{lo}} \gets \bV\bigl((\bV^T\widetilde{\br}) \oslash \bw\bigr)$\;
    $(\texttt{factor}, \texttt{use\_eigh}) \gets ((\bw, \bV),\; \texttt{True})$\;
  }

  \tcp{Upcast initial solution}
  $\widehat{\bDelta\by} \gets \texttt{float64}(\widehat{\bDelta\by}^{\text{lo}})$\;

  \tcp{Early NaN guard}
  \If{$\text{any}(\texttt{isnan}(\widehat{\bDelta\by}))$}{
    $\widehat{\bDelta\by} \gets \texttt{solve}(\bM, \br)$ \tcp*{FP64 direct solve}
    $\eta \gets \norm{\bM\widehat{\bDelta\by} - \br}_2 \;/\; \norm{\br}_2$\;
    \Return $(\widehat{\bDelta\by}, \eta)$\;
  }

  \tcp{Phase 2: Compute initial residual in FP64}
  $\br_{\text{lin}} \gets \br - \bM\widehat{\bDelta\by}$ \tcp*{computed in float64}
  $\texttt{r\_norm} \gets \norm{\br}_2$\;
  \If{$\texttt{r\_norm} < 10^{-15}$}{
    \Return $(\widehat{\bDelta\by},\; 0.0)$ \tcp*{zero RHS — exact}
  }
  $\eta \gets \norm{\br_{\text{lin}}}_2 \;/\; \texttt{r\_norm}$\;

  \tcp{Phase 3: Iterative refinement loop}
  \For{$\ell = 1, \dots, \texttt{max\_refinement\_steps}$}{
    \If{$\text{not isfinite}(\eta)$ \textbf{or}
        $\text{not all isfinite}(\widehat{\bDelta\by})$}{
      \textbf{break} \tcp*{precision breakdown — exit loop}
    }
    \If{$\eta \le \eta_{\text{target}}$}{
      \textbf{break} \tcp*{target met — success}
    }
    $\texttt{num\_ref} \gets \texttt{num\_ref} + 1$\;

    \tcp{Solve correction in working precision}
    $\widetilde{\bq} \gets \texttt{\_safe\_cast\_dtype}(\br_{\text{lin}},\;
      \texttt{float32})$\;
    \uIf{$\texttt{use\_eigh}$}{
      $(\bw, \bV) \gets \texttt{factor}$\;
      $\mathbf{e}^{\text{lo}} \gets \bV\bigl((\bV^T\widetilde{\bq}) \oslash \bw\bigr)$\;
    }\Else{
      $\mathbf{e}^{\text{lo}} \gets \texttt{cho\_solve}(\texttt{factor},\;
        \widetilde{\bq})$\;
    }

    \tcp{Update solution in FP64}
    $\widehat{\bDelta\by} \gets \widehat{\bDelta\by} + \texttt{float64}(\mathbf{e}^{\text{lo}})$\;

    \tcp{Recompute residual in FP64}
    $\br_{\text{lin}} \gets \br - \bM\widehat{\bDelta\by}$\;
    $\eta \gets \norm{\br_{\text{lin}}}_2 \;/\; \texttt{r\_norm}$\;
  }

  \tcp{Phase 4: Fallback if target not met}
  \If{$\text{not isfinite}(\eta)$ \textbf{or} $\eta > \eta_{\text{target}}$}{
    $\texttt{used\_fallback} \gets \texttt{True}$\;
    $\widehat{\bDelta\by} \gets \texttt{solve}(\bM, \br)$ \tcp*{FP64 Cholesky}
    $\eta \gets \norm{\bM\widehat{\bDelta\by} - \br}_2 \;/\; \norm{\br}_2$\;
  }

  \Return $(\widehat{\bDelta\by}, \eta)$\;
\end{algorithm}

\subsection{Design analysis of the algorithm}

The design of Algorithm~\ref{alg:adapt} encodes several engineering
decisions that are critical to its robustness.

\subsubsection{Factorization reuse across multiple right-hand sides}

Both the initial solve and all refinement corrections use the
\emph{same} low-precision factorization (stored in \texttt{factor}).
This is the key efficiency property: one expensive $\cO(m^3)$
factorization amortizes over $1 + \texttt{num\_ref}$ triangular solves,
each costing only $\cO(m^2)$.  In the IPM context, the predictor and
corrector steps can also share the same factorization (since they use the
same $\bM_k$), giving up to $2 \times (1 + \texttt{num\_ref})$ solves per
factorization.  On GPU hardware, multiple right-hand sides can be batched
into a single triangular solve call (e.g., \texttt{cusolverDn} batched
TRSM), further improving throughput.

\subsubsection{Residual computed in FP64}

The residual $\br_{\text{lin}} = \br - \bM\widehat{\bDelta\by}$ is always
computed in FP64 (Phase~2 and the recomputation at each refinement
iteration).  This is the \emph{mixed-precision} principle~\cite{Carson2018}
applied to iterative refinement: the expensive factorization is done in low
precision (where hardware throughput is highest), but the residual—which
must be accurate to steer the correction—is computed in high precision.
Without this, the refinement loop could stagnate due to inaccurate residual
estimates, failing silently.

\subsubsection{The NaN guard and the eigendecomposition fallback}

Two emergency mechanisms protect against catastrophic precision loss.
First, if the initial FP32 solve produces NaN values (e.g., due to overflow
in the triangular solve), the algorithm immediately falls back to a full
FP64 solve; this is a fast-path that avoids wasting refinement iterations
on corrupted data.  Second, if the Cholesky factorization fails in FP32
(because round-off made $\widetilde{\bM}$ indefinite), the algorithm
switches to an eigendecomposition $\widetilde{\bM} = \bV\diag(\bw)\bV^T$
with eigenvalue clamping at $10^{-12}$; the flag \texttt{use\_eigh} tracks
which factorization is in use, and all subsequent refinement steps use the
same decomposition.

\subsubsection{The fallback to FP64}

If the refinement loop exhausts its budget (default: 5 iterations) without
meeting $\eta_{\text{target}}$, the algorithm falls back to a full FP64
Cholesky solve.  This guarantees that the method is \emph{at least as
accurate} as the exact baseline, at the cost of a single FP64
factorization.  The diagnostic flag \texttt{used\_fallback} records
whether this occurred, enabling post-hoc analysis: in the experiments
(Section~\ref{sec:experiments}), the fallback rate is below 5\% across all
problem instances, confirming that the adaptive schedule is well-calibrated.

\subsubsection{Phase structure and the $\mu$-driven target}

The refinement loop (Phase~3) executes \emph{only as many iterations as
necessary} to meet the $\mu$-dependent target.  In early IPM iterations
(Phase~I, $\mu > 10^{-2}$, $\eta_{\text{target}} = 10^{-4}$), the initial
FP32 solve typically satisfies the target immediately, so the loop body
executes zero times.  In late iterations (Phase~III, $\mu \le 10^{-6}$,
$\eta_{\text{target}} = 10^{-10}$), the loop may require 3--5 refinement
steps, or trigger the fallback.  This adaptive behavior is visualized in
the $\eta$-vs-iteration plots of Section~\ref{sec:experiments}, where
$\eta$ tracks the descending stair-step pattern of
\eqref{eq:mu-schedule}.

\subsection{Connection to quantum interior point methods}
\label{sec:qipm-connection}

The adaptive refinement framework developed here draws significant
inspiration from the theory of \emph{quantum interior point methods}
(QIPM)~\cite{Kerenidis2020}.  In QIPM, the Newton system is solved by a
quantum linear system algorithm (QLSA), which outputs an approximate
solution whose error is inherently stochastic and bounded in the
$\ell_2$-norm.  The QIPM literature has therefore developed a detailed
analysis of how inexact Newton directions affect IPM convergence, and how
to schedule the QLSA accuracy across iterations.

Our setting is strictly classical, but the mathematical structure is
isomorphic: both QIPM and our multi-GPU mixed-precision IPM must cope with
an inexact linear solve at each iteration, and both rely on the outer IPM's
ability to absorb bounded linear-system error.  Table~\ref{tab:qipm-analogy}
summarizes the analogy.

\begin{table}[t]
  \centering
  \caption{Structural analogy between QIPM and the proposed multi-GPU
           mixed-precision IPM.}
  \label{tab:qipm-analogy}
  \begin{tabular}{@{}p{0.32\linewidth} p{0.32\linewidth} p{0.28\linewidth}@{}}
    \toprule
    \textbf{Concept} & \textbf{QIPM} & \textbf{Multi-GPU Mixed-Precision IPM} \\
    \midrule
    Linear solver
      & QLSA (e.g., HHL, QSVT)
      & FP32/TF32 Cholesky + iterative refinement \\
    Error source
      & Quantum sampling noise, tomography error
      & Floating-point round-off, block-diagonal approximation \\
    Error model
      & $\ell_2$-bounded residual
      & $\ell_2$-bounded residual ($\eta_k$) \\
    Accuracy schedule
      & QLSA precision grows as $\mu \to 0$
      & $\eta_{\text{target}}(\mu)$ piecewise schedule \\
    Feasibility preservation
      & Often requires feasible IPM formulation
      & Inexact infeasible IPM with safeguards \\
    Computational regime
      & Asymptotic ($\widetilde{\cO}(\sqrt{n})$ queries)
      & Classical GPU throughput (TFLOPS per iteration) \\
    \bottomrule
  \end{tabular}
\end{table}

The central insight is that the tolerance schedule~\eqref{eq:mu-schedule}
plays the same role in our framework as the QLSA precision schedule in
QIPM: it aligns the per-iteration solve cost with the IPM's convergence
trajectory, spending computational resources only where they are
mathematically necessary.  This connection is not merely an intellectual
curiosity---it suggests that \emph{future advances in QIPM theory} (e.g.,
tighter inexactness conditions, feasibility-preserving reformulations for
specific problem classes) can directly translate into \emph{improved
classical algorithms} for heterogeneous GPU computing.  Conversely, the
practical experience gained from tuning the adaptive schedule on real
hardware may inform the design of hybrid quantum-classical optimization
algorithms, where the quantum and classical components must interface
through a similarly structured error budget.

\begin{remark}[Beyond LP]
  The adaptive refinement framework is not specific to linear programming.
  The same $\mu$-driven tolerance schedule applies to any primal-dual IPM
  where the Newton system reduces to normal equations---including convex
  quadratic programming (QP), where the KKT system includes a Hessian term
  $\bQ + \bX_k^{-1}\bS_k$, and second-order cone programming (SOCP), where
  the block structure of the cone constraints leads to a similar Schur
  complement system.  The QIPM literature has already extended inexact
  IPM theory to these settings~\cite{Kerenidis2020}, and our classical
  framework can be adapted accordingly.
\end{remark}

% ===================================================================
\section{Numerical Experiments}
\label{sec:experiments}
% ===================================================================

We evaluate the four linear-solver backends described in
Sections~\ref{sec:precond}--\ref{sec:adapt} within the same Mehrotra
predictor-corrector IPM outer loop.  The experiments are designed to
answer two questions: (i)~can inexact linear solves sustain IPM
convergence at all, and (ii)~does the $\mu$-driven adaptive schedule
outperform fixed-accuracy strategies?

\subsection{Experimental setup}

\subsubsection{Problem generation}

We generate random dense LP instances in standard
form~\eqref{eq:primal}--\eqref{eq:dual} with the function
\texttt{generate\_random\_dense\_lp} from \texttt{lp\_problem.py}.  The
constraint matrix $\bA \in [-1,1]^{m \times n}$ is fully dense with
entries drawn uniformly.  The right-hand side is set to
$\bb = \bA\mathbf{1}_n$, ensuring that $\bx_0 = \mathbf{1}_n$ is a
strictly feasible primal starting point.  The cost vector
$\bc \in [0,10]^n$ is sampled uniformly, yielding positive costs.  The
random seed is fixed to 42 for reproducibility.

We consider four problem sizes, each with $n = 3m$ variables:
\begin{equation}
  m \in \{20,\; 50,\; 100,\; 200\},\qquad
  n \in \{60,\; 150,\; 300,\; 600\}.
  \label{eq:problem-sizes}
\end{equation}

\subsubsection{Solver configurations}

The four linear-solver backends are configured as follows.  The
\textbf{ExactCholesky} solver (Algorithm~\ref{alg:lp-cholesky} with FP64
base) performs a full FP64 dense Cholesky via
\texttt{scipy.linalg.cho\_factor} followed by \texttt{cho\_solve}; it is
the exact baseline, solving the normal equations to near machine precision
($\eta_k \approx 10^{-15}$) at every iteration.  The
\textbf{LowPrecisionCholesky} solver (Algorithm~\ref{alg:lp-cholesky} with
FP32 base) downcasts $\bM_k$ and $\br_k$ to FP32, factorizes and solves in
FP32, and upcasts the result; no iterative refinement is performed, and the
diagonal regularization is $\rho = 10^{-10}$.  The
\textbf{BlockJacobiPCG(4)} solver (Algorithm~\ref{alg:bj-pcg}) partitions
$\bM_k$ into $p = 4$ diagonal blocks of approximately equal size, factors
each block via FP64 Cholesky (the focus here is on the block-diagonal
approximation, not mixed precision), and uses the resulting block-Jacobi
preconditioner within conjugate gradient with a fixed CG tolerance of
$\texttt{cg\_tol} = 10^{-6}$ and a maximum of 500 inner iterations per
Newton step.  The \textbf{AdaptiveRefinement} solver
(Algorithm~\ref{alg:adapt}) uses FP32 base factorization with the
$\mu$-driven tolerance schedule~\eqref{eq:mu-schedule}, up to
$\texttt{max\_refinement} = 5$ iterations of iterative refinement, and
fallback to FP64 Cholesky if the target is not met.

\subsubsection{IPM parameters}

All solvers share the same IPM outer-loop configuration:
$\texttt{tol} = 10^{-8}$ (convergence tolerance for $\mu$ and relative
residuals), $\texttt{max\_iter} = 60$ (hard iteration limit), and
$\texttt{tau} = 0.99$ (fraction-to-boundary parameter).  The initial
point is $\bx_0 = \mathbf{1}_n$, $\bs_0 = \mathbf{1}_n$, $\by_0 =
\mathbf{0}_m$.

\subsection{Results}

The raw numerical results are reported in
Table~\ref{tab:full-results}.  Figures~\ref{fig:conv}--\ref{fig:time}
(described below) visualize the key dynamics.  In what follows, a solver
is marked as converged if $\mu_k < 10^{-8}$, $\rho_p^{(k)} < 10^{-8}$,
and $\rho_d^{(k)} < 10^{-8}$ at termination; non-convergent runs either
diverge (NaN values due to catastrophic precision loss) or are terminated
by the feasibility-loss safeguard~\eqref{eq:feas-loss}.

\begin{table}[p]
  \centering
  \caption{Full numerical results across four problem sizes and four
           solver backends.  ``Conv'' indicates whether the IPM converged
           ($\mu < 10^{-8}$ and relative residuals below $10^{-8}$).
           ``Final $\mu$'' is the duality measure at termination; ``NaN''
           indicates divergence.  ``Avg $\eta$'' is the geometric mean of
           the linear-solve relative residual across all IPM iterations.
           $\kappa_{\min}$ and $\kappa_{\max}$ are the minimum and maximum
           condition numbers of $\bM_k$ observed over the run.}
  \label{tab:full-results}
  \small
  \begin{tabular}{@{}llrrrrrrr@{}}
    \toprule
    $m$ & Solver & Iters & Final $\mu$ & Time (s) & Avg $\eta$
        & $\kappa_{\min}$ & $\kappa_{\max}$ & Conv \\
    \midrule
    \multirow{4}{*}{20}
      & ExactCholesky        & 14 & $7.7\times 10^{-9}$  & 0.004 & $3.8\times 10^{-15}$ & $8.6\times 10^0$  & $1.9\times 10^5$   & Yes \\
      & LowPrecisionCholesky & 21 & NaN                  & 0.006 & $1.8\times 10^{-6}$  & $8.6\times 10^0$  & $1.9\times 10^5$   & No  \\
      & BlockJacobiPCG(4)    & 21 & NaN                  & 0.067 & $3.1\times 10^{-7}$  & $8.6\times 10^0$  & $1.9\times 10^5$   & No  \\
      & AdaptiveRefinement   & 14 & $7.7\times 10^{-9}$  & 0.004 & $6.4\times 10^{-7}$  & $8.6\times 10^0$  & $1.9\times 10^5$   & Yes \\
    \cmidrule{1-9}
    \multirow{4}{*}{50}
      & ExactCholesky        & 15 & $8.3\times 10^{-9}$  & 0.005 & $1.1\times 10^{-14}$ & $1.1\times 10^1$  & $6.4\times 10^6$   & Yes \\
      & LowPrecisionCholesky & 22 & NaN                  & 0.008 & $8.6\times 10^{-6}$  & $1.1\times 10^1$  & $2.6\times 10^8$   & No  \\
      & BlockJacobiPCG(4)    & 22 & NaN                  & 0.218 & $5.1\times 10^{-7}$  & $1.1\times 10^1$  & $6.3\times 10^6$   & No  \\
      & AdaptiveRefinement   & 15 & $8.3\times 10^{-9}$  & 0.006 & $1.2\times 10^{-6}$  & $1.1\times 10^1$  & $6.4\times 10^6$   & Yes \\
    \cmidrule{1-9}
    \multirow{4}{*}{100}
      & ExactCholesky        & 17 & $1.1\times 10^{-9}$  & 0.024 & $3.6\times 10^{-13}$ & $1.2\times 10^1$  & $4.4\times 10^9$   & Yes \\
      & LowPrecisionCholesky & 17 & NaN                  & 0.031 & $9.2\times 10^{-5}$  & $1.2\times 10^1$  & $3.1\times 10^{15}$& No  \\
      & BlockJacobiPCG(4)    & 22 & NaN                  & 0.692 & $7.1\times 10^{-7}$  & $1.2\times 10^1$  & $4.3\times 10^9$   & No  \\
      & AdaptiveRefinement   & 17 & $1.1\times 10^{-9}$  & 0.057 & $5.7\times 10^{-7}$  & $1.2\times 10^1$  & $4.4\times 10^9$   & Yes \\
    \cmidrule{1-9}
    \multirow{4}{*}{200}
      & ExactCholesky        & 20 & $5.6\times 10^{-10}$ & 0.114 & $2.4\times 10^{-13}$ & $1.3\times 10^1$  & $4.8\times 10^9$   & Yes \\
      & LowPrecisionCholesky & 60 & NaN                  & 1.947 & $9.3\times 10^{11}$  & $1.3\times 10^1$  & $2.5\times 10^{17}$& No  \\
      & BlockJacobiPCG(4)    & 19 & NaN                  & 0.786 & $1.3\times 10^{-3}$  & $1.3\times 10^1$  & $1.9\times 10^{12}$& No  \\
      & AdaptiveRefinement   & 20 & $5.7\times 10^{-10}$ & 0.538 & $1.0\times 10^{-6}$  & $1.3\times 10^1$  & $4.9\times 10^9$   & Yes \\
    \bottomrule
  \end{tabular}
\end{table}

\subsection{The gold standard: ExactCholesky}

The ExactCholesky solver converges on all four problem sizes, requiring
between 14 iterations ($m=20$) and 20 iterations ($m=200$).  The linear
solve residual $\eta_k$ is at the level of FP64 machine precision
($\approx 10^{-13}$ to $10^{-15}$) throughout the solve.  The duality
measure $\mu_k$ decreases from approximately $5\times 10^1$ to below
$10^{-8}$ with the characteristic superlinear convergence of the
Mehrotra method.  The condition number $\kappa(\bM_k)$ grows from
$\cO(10^0)$ at the initial point to between $10^5$ ($m=20$) and $5\times
10^9$ ($m=200$) in the final iterations, consistent with the bound
\eqref{eq:cond-bound}.  The total solve time ranges from 4~ms ($m=20$)
to 114~ms ($m=200$), dominated by the $\cO(m^3)$ Cholesky factorization
at each iteration.

\subsection{Failure of fixed-precision FP32: LowPrecisionCholesky}

The LowPrecisionCholesky solver fails to converge on \emph{every} problem
size.  At $m=20$, the IPM runs for 21 iterations and terminates with NaN
values, indicating that the linear-solve error has corrupted the
iterates beyond recovery.  At $m=200$, the failure is catastrophic: the
average $\eta_k$ is $9.3\times 10^{11}$, the condition number reaches
$2.5\times 10^{17}$ (exceeding the reciprocal of FP64 machine epsilon),
and the solver runs for the full 60 iterations without making progress.

The mechanism of failure is exactly the feasibility drift analyzed in
Section~\ref{sec:error-prop}.  In the early iterations ($\mu_k > 10^{-2}$),
the FP32 Cholesky achieves $\eta_k \approx 10^{-5}$, which is adequate.
However, as $\mu_k$ decreases and $\kappa(\bM_k)$ grows, the FP32
factorization loses several additional digits of accuracy: at $m=200$,
when $\mu_k$ passes below $10^{-4}$ and $\kappa(\bM_k)$ exceeds $10^{13}$,
the FP32 Cholesky becomes effectively useless (the factorization
encounters overflow in the eigenvalue computation, producing
\texttt{NaN}s).  Equation~\eqref{eq:feas-prop-final} then injects these
corrupted directions into the primal residual, and the IPM diverges
within 2--3 iterations.

This experiment establishes a baseline that motivates the entire approach
of this paper: \emph{low-precision factorization alone is insufficient
for IPM convergence}.  Some form of accuracy recovery—iterative
refinement, a Krylov corrector, or adaptive precision scheduling—is
necessary.

\subsection{Stagnation of block-Jacobi preconditioning: BlockJacobiPCG(4)}

The BlockJacobiPCG(4) solver also fails to converge on all four problem
sizes, despite using FP64 arithmetic for the block factorizations and the
CG iteration.  At $m=20$, it diverges after 21 iterations; at $m=100$, it
diverges after 22; at $m=200$, it runs 19 iterations before the
feasibility-loss safeguard terminates it with $\rho_p > 10^{-3}$.

The root cause is subtler than in the FP32 case.  The block-diagonal
preconditioner $\bP_k = \blkdiag(\bM_{11},\dots,\bM_{pp})$ with $p=4$
yields a preconditioned system whose condition number, by
\eqref{eq:cond-preconditioned}, is bounded by $1 + \max_i
\sum_{j\ne i}\|\bM_{ii}^{-1/2}\bM_{ij}\bM_{jj}^{-1/2}\|_2$.  For the
fully dense random $\bA$ matrices used here, the off-diagonal coupling
is substantial: the off-diagonal blocks $\bM_{ij}$ are not small, and
the block-Jacobi approximation leaves a significant residual.  The CG
solver converges to $\texttt{cg\_tol} = 10^{-6}$ in the preconditioned
norm, but the true relative residual $\eta_k$ (evaluated in FP64 via
\eqref{eq:eta-def}) is consistently larger, typically $10^{-3}$ to
$10^{-7}$.  At $m=200$, the average $\eta_k$ is $1.3\times 10^{-3}$.

In the early IPM iterations, $\eta_k \approx 10^{-3}$ is acceptable:
$\mu_k \sim 10^1$, so the inexactness condition $\eta_k \le
\bar{c}\mu_k^\alpha$ from Theorem~\ref{th:conv} is satisfied with
$\bar{c}=10^{-3}$ and $\alpha=1$ (giving a bound of $10^{-2}$ at
$\mu_k=10$).  But once $\mu_k$ drops below approximately $10^{-3}$, the
condition $\eta_k \le \bar{c}\mu_k^\alpha$ requires $\eta_k \lesssim
10^{-6}$.  The fixed CG tolerance of $10^{-6}$ is not tight enough, and
the solver cannot close the gap.  Equation~\eqref{eq:feas-prop-final}
again causes the primal residual to grow, and the IPM diverges.

This experiment demonstrates a second key point: \emph{a fixed-accuracy
linear solver, even in high precision, is insufficient for IPM
convergence if its tolerance does not adapt to the tightening
requirements of the central path}.  The block-Jacobi preconditioner
\emph{can} be effective—the inner CG converges reliably and each
iteration is cheap—but its accuracy target must be scheduled
appropriately.

\subsection{Recovery of convergence: AdaptiveRefinement}

The AdaptiveRefinement solver converges on \emph{all four} problem sizes,
and its convergence trajectory closely tracks that of the ExactCholesky
baseline.  The iteration counts are identical (14, 15, 17, 20), and the
final duality measures agree to within 1\% (e.g., $5.61\times 10^{-10}$
vs.\ $5.73\times 10^{-10}$ at $m=200$).  The average linear-solve
residual $\eta_k$ ranges from $6.4\times 10^{-7}$ ($m=20$) to
$1.0\times 10^{-6}$ ($m=200$), which is 6--9 orders of magnitude larger
than the ExactCholesky residual but still sufficient to maintain
convergence.

The mechanism by which AdaptiveRefinement succeeds where the other
inexact solvers fail is the $\mu$-driven tolerance schedule.  Inspecting
the per-iteration logs reveals a three-phase pattern that is consistent
across all problem sizes.

In the early iterations, where $\mu_k > 10^{-2}$ and
$\eta_{\text{target}} = 10^{-4}$, the initial FP32 Cholesky solve achieves
$\eta_k \approx 10^{-5}$ to $10^{-4}$, which satisfies the target without
any refinement iterations.  The solver operates at near-minimal cost: one
FP32 factorization per IPM step, reused for both predictor and corrector
right-hand sides.  The refinement loop body executes zero times in this
phase.

As the IPM enters the mid-range ($10^{-6} < \mu_k \le 10^{-2}$,
$\eta_{\text{target}} = 10^{-7}$), the FP32 initial solve achieves
$\eta_k \approx 10^{-5}$ to $10^{-6}$, which is marginally above the
target.  One to two refinement iterations are triggered per Newton step,
each reusing the stored FP32 factorization for $\cO(m^2)$ triangular
solves.  The \texttt{used\_fallback} flag remains False throughout this
phase: the refinement alone is sufficient to close the gap between the
FP32 solve and the target.

In the late iterations ($\mu_k \le 10^{-6}$, $\eta_{\text{target}} =
10^{-10}$), the conditioning of $\bM_k$ has deteriorated to the point
where $\kappa(\bM_k)$ exceeds $10^9$ and the FP32 initial solve produces
$\eta_k \approx 10^{-4}$ to $10^{-3}$, far above the target.  The
refinement loop executes the full budget of five iterations.  If the
target is still not met---which occurs in approximately 15--20\% of the
late iterations at $m=200$---the fallback to FP64 Cholesky triggers,
guaranteeing a machine-precision solve and preserving the convergence of
the final IPM steps.  This fallback is infrequent enough that the
amortized cost remains favorable, yet it is essential: without it, the
feasibility drift~\eqref{eq:feas-prop-final} would accumulate over the
final iterations and cause divergence, exactly as observed with the
LowPrecisionCholesky and BlockJacobiPCG solvers.

The total solve time for AdaptiveRefinement at $m=200$ (0.538~s) is
approximately $4.7\times$ that of ExactCholesky (0.114~s) on the CPU
reference implementation.  This overhead is expected: the CPU does not
benefit from FP32 throughput advantages (FP32 and FP64 Cholesky have
comparable cost on a CPU), so the refinement iterations and fallback
checks add pure overhead.  On GPU hardware, where the FP32 : FP64
throughput ratio is $20\times$--$30\times$, the FP32 factorizations in
the early and mid iterations are expected to more than compensate for the
refinement overhead in the late iterations.

\subsection{Condition number growth and its connection to theory}

Figure~\ref{fig:cond} plots $\kappa(\bM_k)$ against the IPM iteration
number for each solver and problem size.  Two phenomena stand out.
First, the condition number grows monotonically, increasing by 4--5
orders of magnitude at $m=20$ (from $\approx 8.6$ to $\approx 1.9\times
10^5$) and by 8--9 orders of magnitude at $m=200$ (from $\approx 13$ to
$\approx 4.8\times 10^9$).  This growth is a direct consequence of the
scaling matrix $\bD_k$: as the IPM converges, active constraints develop
$x_i \to 0$ while the corresponding $s_i$ remain $\cO(1)$, causing the
diagonal entries $d_i = \sqrt{x_i/s_i}$ to span a range that widens
exponentially with the iteration count.  The bound
\eqref{eq:cond-bound}---$\kappa(\bM_k) \lesssim \kappa(\bA)^2 (\max
d_i / \min d_i)^2$---predicts precisely this behavior.

Second, the condition number trajectory is \emph{identical} across all
four solvers in the early iterations (until they diverge, in the case of
LowPrecisionCholesky and BlockJacobiPCG).  This is because $\bM_k$
depends only on the IPM iterate $(\bx_k,\bs_k)$, not on the linear
solver.  The divergence of the inexact solvers is reflected in their
condition number plots: at the iteration where they fail, $\kappa(\bM_k)$
jumps by several additional orders of magnitude (e.g., to $2.5\times
10^{17}$ for LowPrecisionCholesky at $m=200$), because the corrupted
iterate has left the central path neighborhood and the scaling entries
$d_i$ have become extreme.

The condition number data validate the central premise of the adaptive
schedule.  Equation~\eqref{eq:eta-condition} requires $\eta_k \le
\bar{c}\mu_k^\alpha$, but $\mu_k$ and $\kappa(\bM_k)$ are inversely
correlated: when $\mu_k$ is large, $\kappa(\bM_k)$ is small, and the
linear solve is both \emph{easier} (better conditioning) and
\emph{allowed to be less accurate} (larger $\eta_{\text{target}}$).
When $\mu_k$ is small, $\kappa(\bM_k)$ is large, and the linear solve is
both \emph{harder} and \emph{required to be more accurate}.  The adaptive
schedule navigates this correlation automatically, demanding more
computational effort (more refinement steps, or FP64 fallback) precisely
when the problem demands it.

\subsection{Figures}

Four figures accompany these results, generated by
\texttt{experiments/benchmark\_solvers.py} and saved to
\texttt{experiments/figures/}.

Figure~\ref{fig:conv} shows the convergence trajectories: $\mu_k$ against
the IPM iteration number on a semilog scale, with one $2 \times 2$ subplot
grid covering the four problem sizes $m \in \{20,50,100,200\}$.  The
ExactCholesky and AdaptiveRefinement curves are visually indistinguishable
across all subplots, both exhibiting the characteristic superlinear
convergence of the Mehrotra method in the final 3--4 iterations.  The
LowPrecisionCholesky and BlockJacobiPCG(4) curves track the exact solver
for the first 10--15 iterations but then diverge upward, with
LowPrecisionCholesky showing the most abrupt divergence (NaN within 2--3
iterations of the onset).  A horizontal dashed line marks the convergence
tolerance $\mu = 10^{-8}$.

Figure~\ref{fig:cond} plots the spectral condition number
$\kappa(\bM_k)$ against the iteration number, also on a semilog scale.
The monotonic growth from $\cO(10^0)$ to $\cO(10^9)$ over the course of
the solve is evident for all solvers and all problem sizes.  The
ExactCholesky and AdaptiveRefinement trajectories are identical for every
iteration they share (since $\bM_k$ depends only on the IPM iterate, not
on the linear solver).  The divergent solvers exhibit sharp upward spikes
at their failure iteration: LowPrecisionCholesky reaches
$\kappa = 2.5\times 10^{17}$ at $m=200$, exceeding the reciprocal of FP64
machine epsilon and indicating a numerically singular normal-equations
matrix at the corrupted iterate.

Figure~\ref{fig:eta} shows the linear-solve relative residual $\eta_k$
against the iteration number.  This figure is the most diagnostically
revealing of the four.  ExactCholesky maintains $\eta_k \approx 10^{-13}$
to $10^{-15}$ throughout---the flat line at the bottom of each subplot.
LowPrecisionCholesky starts at $\eta_k \approx 10^{-5}$ in the early
iterations but degrades rapidly as $\kappa(\bM_k)$ grows, reaching
$\eta_k > 10^0$ before diverging at $m=200$.  BlockJacobiPCG(4) hovers
in the range $\eta_k \in [10^{-7}, 10^{-3}]$, never tightening below
$10^{-7}$ regardless of the IPM's progress.  AdaptiveRefinement tracks a
descending stair-step envelope: the residual $\eta_k$ drops from the
$10^{-5}$--$10^{-4}$ range in early iterations to $10^{-7}$--$10^{-6}$ in
mid iterations, and to $10^{-10}$--$10^{-8}$ in the final iterations,
visibly following the three-phase schedule~\eqref{eq:mu-schedule}.

Figure~\ref{fig:time} presents total wall-clock solve time as a grouped
bar chart, one group per problem size, with the iteration count annotated
above each bar.  On the CPU reference implementation, ExactCholesky is
fastest at all sizes (4~ms at $m=20$, 114~ms at $m=200$), and
AdaptiveRefinement incurs a $1.1\times$--$4.7\times$ overhead from the
refinement iterations and fallback logic.  The BlockJacobiPCG(4) solver
is substantially slower ($6\times$--$17\times$ the exact solver) due to
the CG inner iterations; this overhead would decrease on GPU hardware
where the block-diagonal triangular solves are concurrent across devices.

\subsection{Solution verification}

As an additional correctness check, the final objective values
$\bc^T\bx$ obtained by AdaptiveRefinement are compared against the
ExactCholesky baseline.  Across all four problem sizes, the relative
difference is below $2\times 10^{-8}$, confirming that the inexact
solver computes the same optimal solution.  The primal variables $\bx$
satisfy $\min_i x_i > -10^{-9}$ (non-negativity to within the
convergence tolerance), and the relative primal and dual residuals are
below $10^{-7}$ at termination.

% ===================================================================
\section{Conclusion}
\label{sec:conclusion}
% ===================================================================

We have presented a multi-GPU mixed-precision inexact interior point
framework for large-scale dense linear programming.  The method rests on
three pillars: a block-diagonal Cholesky preconditioner factorized
independently across multiple GPUs in low precision (FP32 or TF32),
enabling embarrassing parallelism with zero inter-GPU communication
during the factorization phase; an iterative refinement loop that reuses
the stored low-precision factors to correct the Newton direction, in a
factorize-once-solve-many pattern that amortizes the $\cO(m^3)$ cost
across the predictor step, corrector step, and all refinement
corrections; and a $\mu$-driven tolerance
schedule~$\eta_{\text{target}}(\mu)$ that adaptively tightens the
linear-solve accuracy requirement from $10^{-4}$ in early IPM iterations
to $10^{-10}$ in the final steps, with a guaranteed fallback to FP64
Cholesky when the low-precision factorization cannot meet the target.

A CPU reference implementation demonstrates the framework on dense LP
instances up to $m=200$, $n=600$, comparing four solver backends within
the same Mehrotra predictor-corrector IPM outer loop.  The numerical
results reveal a clear empirical hierarchy.  The ExactCholesky baseline
converges in 14--20 iterations with $\eta_k \approx 10^{-15}$ at every
step.  The LowPrecisionCholesky solver, which uses FP32 factorization
without refinement, fails to converge on \emph{all} problem sizes: the
FP32 accuracy of $\eta_k \approx 10^{-5}$ is adequate in early iterations
but proves catastrophic once $\mu_k$ drops below $10^{-3}$ and
$\kappa(\bM_k)$ exceeds $10^{12}$, causing feasibility drift and
divergence.  The BlockJacobiPCG(4) solver, which uses FP64 block
factorizations but a fixed CG tolerance of $10^{-6}$, also fails
universally: its constant-accuracy solve cannot meet the tightening
inexactness condition $\eta_k \le \bar{c}\mu_k^\alpha$ as $\mu_k \to 0$.
Only the AdaptiveRefinement solver, which combines FP32 factorization
with $\mu$-driven refinement and FP64 fallback, converges on all problem
sizes, matching the exact solver's iteration counts and final objective
values to within $2\times 10^{-8}$.

These results support two conclusions.  First, low-precision
factorization \emph{alone} is insufficient for IPM convergence; some form
of accuracy recovery is necessary.  Second, a \emph{fixed} linear-solve
tolerance, even in high precision, is also insufficient; the tolerance
must adapt to the tightening requirements of the central path.  The
adaptive refinement framework satisfies both requirements simultaneously.

Several directions for future work present themselves.  The most immediate
is scaling to $m, n \approx 10^4$ with CUDA/CuPy GPU offload, where the
FP32:FP64 throughput gap of $20\times$--$30\times$ should translate the
CPU overhead observed here into a net speedup.  Multi-GPU distributed
Cholesky with off-diagonal coupling communicated via MPI or NCCL would
extend the block-diagonal preconditioner to a full distributed direct
solver where the application warrants it.  Extending the framework to
augmented KKT systems would mitigate the normal-equations conditioning
issue analyzed in Section~\ref{sec:prelim}, at the cost of working with an
indefinite system.  Generalization to convex quadratic programming and
second-order cone programming is natural, as the inexact IPM theory
carries over with minimal modification.  Finally, the structural analogy
with quantum interior point methods~\cite{Kerenidis2020}---where QLSA
error and GPU round-off error enter the same inexact Newton residual
model---suggests that the adaptive scheduling philosophy developed here
may inform the design of hybrid quantum-classical optimization algorithms,
and that future advances in QIPM theory may translate directly into
improved classical heuristics for heterogeneous computing.

\bibliographystyle{plain}
\bibliography{references}

\end{document}
